\BourbakiExerciseGroup

除非另有明确说明，本文所考虑的所有环均假定为交换环。

\begin{enumerate}
\def\labelenumi{\arabic{enumi})}
\tightlist
\item
  设 A 是一个全序环，B 是 A 的一个加法子群，且 B 对乘法封闭。
\end{enumerate}

\begin{enumerate}
\def\labelenumi{\alph{enumi})}
\item
  元素 \(x \in A\) 称为关于 B 是无穷大的，如果 \(|y| < |x|\) 对所有 \(y \in B\) 成立。证明 A 中关于 B 不是无穷大的元素所成的集 \(F(B)\) 是 A 的一个包含 B 的子环。
\item
  元素 \(x \in A\) 称为关于 B 是无穷小的，如果对 B 中所有 y \textgreater{} 0，都有 \(|x| < y\) 。如果对 B 中所有 y \textgreater{} 0，存在 \(z \in B\) 使得 0 \textless{} z \textless{} y，则 A 中关于 B 为无穷小的元素所成的集是 A 的一个子伪环 I(B)。此外，若对 B 中每一对满足 0 \textless{} y \textless{} z 的元素 y、z，存在 \(x \in B\) 使得 0 \textless{} xz \textless{} y，则 I(B) 是环 F(B) 的一个理想。
\end{enumerate}

\begin{enumerate}
\def\labelenumi{\arabic{enumi})}
\setcounter{enumi}{1}
\item
  设 A 为一个全序环，n 为 A 的幂零元集合（它是 A 的一个理想）。A 中不属于 n 的每个元素关于 n 都是无穷大的；商环 A/n 是全序的（VI，第 30 页，习题 4），并且是一个整环。
\item
  在域 \(\mathbf{Q}\) 中，设 \(\mathbf{P}\) 为由 0 与 \(\geqslant 1\) 的有理数（在通常的序下）所组成的集。证明 \(\mathbf{P}\) 满足公理 \((\mathrm{AP}_{\mathrm{I}})\) 、 \((\mathrm{AP}_{\mathrm{II}})\) 与 \((\mathrm{AP}_{\mathrm{III}})\) 。
\item
  \begin{enumerate}
  \def\labelenumii{\alph{enumii})}
  \tightlist
  \item
    设 K 是一个域，P 是 K 的一个满足条件 \((\mathrm{AP}_{\mathrm{I}})\) 、 \((\mathrm{AP}_{\mathrm{II}})\) 与 \((\mathrm{AP}_{\mathrm{III}})\) 的子集，并且 \(K^{2} \subset P\) （其中 \(K^{2}\) 表示 K 的元素的平方所成的集）；证明在由 P 定义的序中，x \textgreater{} 0 蕴含 \(x^{-1} > 0\) ；由此推出，集合 \(K_{+}^{*}\) （由 K 中 \textgreater{} 0 的元素组成）是 K 的乘法群的一个子群。
  \end{enumerate}
\end{enumerate}

\begin{enumerate}
\def\labelenumi{\alph{enumi})}
\setcounter{enumi}{1}
\item
  在域 \(\mathbf{Q}\) 中，集合 \(\mathbf{P}\) 若满足 \(a\) ) 的条件，则只能是 \(\geqslant 0\) 的有理数（在通常的序下）所成的集。
\item
  * 设 P 是 K 的一个子集，满足 \((\mathrm{AP}_{\mathrm{I}})\) 、 \((\mathrm{AP}_{\mathrm{II}})\) 与 \((\mathrm{AP}_{\mathrm{III}})\) ，使得 \(1 \in P\) ，并且使得在由 P 定义的序中，x \textgreater{} 0 蕴含 \(x^{-1} > 0\) 。证明对正的 y 与 z，\(y^{2} \geqslant z^{2}\) 蕴含 \(y \geqslant z\) 。由此推出，对任意的 y \textgreater{} 0，我们有
\end{enumerate}

\[
\frac {1}{2} (y + y ^ {- 1}) \geqslant 1 - n ^ {- 1}
\]

对每个整数 \(n > 0\) 成立（注意 \((y + y^{-1})^{2m} \geqslant \binom{2m}{m}\) 对每个整数 \(m > 0\) 成立）. 由此推出，若由 P 定义的序加法群结构是阿基米德的（VI，第 35 页，习题 31），则 \((y - z)^2 \geqslant 0\) 对 P 的每一对元素 \(y, z\) 成立；若 \(K'\) 是由 P 生成的 K 的子域，则 \(x^2 \geqslant 0\) 对所有 \(x \in K'\) 成立。 *

\(d) * \text{令 } K = R(X) \text{ 为关于一个不定元的有理函数域，其系数域为 } R; \text{ 令 } P \text{ 为由 0 以及有理函数 } u \in K \text{ 满足以下条件：} u(t) \text{ 有定义且 } >0 \text{ 对每个实数 } t. \text{ 证明 } P \text{ 满足条件 } c) \text{ 并生成 } K, \text{ 但存在元素 } v \in K \text{ 满足以下条件：} v^2 \notin P. *\)

\begin{enumerate}
\def\labelenumi{\arabic{enumi})}
\setcounter{enumi}{4}
\tightlist
\item
  设 A 是一个格序环。A 的一个凸理想（VI，第 30 页，习题 4）称为不可约的，如果它不是两个不同于自身的凸理想的交集。
\end{enumerate}

\begin{enumerate}
\def\labelenumi{\alph{enumi})}
\tightlist
\item
  证明 A 的不可约凸理想的交为 0（若 \(a \in \mathbf{A}\) 非零，考虑 A 的一个在所有不包含 \(a\) 的理想中极大的凸理想）。由此推出 A 同构于一个子环 \(\mathbf{A}'\) （某个乘积 \(\prod_{i} \mathbf{A}_i\) 的），使得每个 \(\mathbf{A}_i\) 都是格序的，
\end{enumerate}

\(\mathrm{pr}_{\iota}(\mathbf{A}^{\prime})=\mathbf{A}_{\iota}\) 对所有 \(\iota\) 成立，且理想 \(\{0\}\) 在 \(\mathbf{A}_{\iota}\) 中（作为凸理想）对所有 \(\iota\) 不可约。 b) 证明以下条件在 A 中等价：

α) \(|xy| = |x| \cdot |y|\) 对所有 x, y；

\(\beta) x \cdot \sup(y, z) = \sup(xy, xz)\) 对所有 \(x \geqslant 0\) , \(y, z\) ;

\(\gamma) x \cdot \inf(y, z) = \inf(xy, xz)\) 对所有 \(x \geqslant 0\) , \(y, z\) ;

\(\delta)\inf (y,z) = 0\) 蕴含 \(\inf (xy,z) = 0\) ，对所有 \(x\geqslant 0\) .

（为看出 \(\gamma\) ) 蕴含 \(\delta\) )，注意 \(xy \leqslant y \cdot \sup (x, 1)\) ).

当这些条件满足时，称A为强格序的。

\begin{enumerate}
\def\labelenumi{\alph{enumi})}
\setcounter{enumi}{2}
\tightlist
\item
  证明一个环 \(\mathbf{A}\) 是强格序的当且仅当它同构于一个乘积 \(\prod_{i} \mathbf{A}_{i}\) 的子环，其中每个因子是全序环。（利用 \(a\) ），归结为证明：若
\end{enumerate}

\{0\} 在强格序环 A 中不可约，则 A 是全序的；为此，注意若 B 是强格序环，则对所有 \(b \in B\) ，集合 \(m\) 由这样的元素 \(x \in B\) 组成：\(\inf(|x|, |b|) = 0\) ；集合 \(n\) 由这样的元素 \(y \in B\) 组成：\(|y| \leqslant |zb|\) 对某个 \(z \in B\) 成立；它们是 B 的两个凸理想，且满足 \(m \cap n = \{0\}\) .

\(d\) ) 证明在强格序环 \(\mathbf{A}\) 中，关系 \(\inf (x,y) = 0\) 蕴含 \(xy = 0\) ；对所有 \(z\in \mathbf{A}\) 有 \(z^2\geqslant 0\) ；对所有 \(x,y\in \mathbf{A}\) 有 \(xy\leqslant \sup (x^2,y^2)\)

\begin{enumerate}
\def\labelenumi{\alph{enumi})}
\setcounter{enumi}{4}
\tightlist
\item
  设 A 是一个没有 \textgreater0 幂零元素的格序环；证明：若 \(\inf(x, y) = 0\) 蕴含 xy = 0，则 A 是强格序的（验证条件 \(\delta\) ) ，即 b) 中所述）。
\end{enumerate}

\begin{enumerate}
\def\labelenumi{\arabic{enumi})}
\setcounter{enumi}{5}
\tightlist
\item
  \begin{enumerate}
  \def\labelenumii{\alph{enumii})}
  \tightlist
  \item
    设 K 是一个格序域。证明以下条件等价：
  \end{enumerate}
\end{enumerate}

\(\alpha) x^{2} \geqslant 0\) 对所有 \(x \in \mathbf{K}\) ;

\(\beta) x > 0\) 蕴含 \(x^{-1} > 0\) ;

γ) K 是强格序的（习题 5）；

δ) K 是全序的。

（为看出 \(\beta\) ) 蕴含 \(\underline{\delta}\) ），注意 \(\beta\) 蕴含 \(xy \leqslant x^2 + y^2\) ，对 \(x, y \in K\) ).

\begin{enumerate}
\def\labelenumi{\alph{enumi})}
\setcounter{enumi}{1}
\item
  设 K 为域 \(\mathbf{Q}(\sqrt{2})\) ，它是由 \(\sqrt{2}\) 添加到有理数域而得到的；作为加法群，K 与 \(Q \times Q\) 通过双射 \(x + y \sqrt{2} \mapsto (x, y)\) 等同；证明：若将 \(Q \times Q\) 上的乘积序（其中 Q 具有通常的序）经由该双射转移到 K 上，则 K 是格序的，但不是强格序的。
\item
  域 \(\mathbf{K} = \mathbf{Q}(\mathbf{X})\) （Q 上关于单个未定元 X 的有理函数域）由它的平方所成的集 \(K^{2}\) 生成。证明：K 中元素的平方和所成的集 P 在 K 上定义了一个与环结构相容的非格序，并且使得 u \textgreater{} 0 蕴含 \(u^{-1} > 0\) .
\end{enumerate}

\begin{enumerate}
\def\labelenumi{\arabic{enumi})}
\setcounter{enumi}{6}
\item
  特征为 \(\neq 2\) 的交换域 K 中，每个元素都是平方和，当且仅当 -1 是平方和（注意 K 的每个元素都是两个平方之差）。
\item
  设 A 是特征为 \(\neq 2\) 的交换域 K 的一个非空子集；若 A 中没有元素是 K 中平方和，则称 K 是 A-可序的。我们称 K 是可序的，如果存在 K 的一个全序，且该全序与其环结构相容。
\end{enumerate}

\begin{enumerate}
\def\labelenumi{\alph{enumi})}
\item
  证明若 K 是 A-可序的，则它是可序的（习题 7），从而特征为 0。
\item
  证明 A-可序域 K 的纯扩张和奇数次代数扩张是 A-可序的（仿照 VI 第 23 页命题 3 论证）。
\item
  设 K 是 A-可序域，b 是 K 中一个元素，且 b 不具有形式 ca - d，其中 \(a \in A\) 而 c 和 d 是 K 中平方和。证明域 \(\mathrm{K}(\sqrt{b})\) 是 A-可序的。
\item
  A-可序域 K 称为极大的，如果 K 没有真代数扩张是 A-可序的。证明极大 A-可序域 K 具有以下性质：1. K 是毕达哥拉斯域，换言之每个平方和都是平方；2. A 中没有元素是平方；3. K 中每个非平方元素都具有形式 \(c^2 a - d^2\) ，其中 \(a \in \mathbf{A}\) ；4. K{[}X{]} 中每个奇数次多项式在 K 中至少有一个根（利用 b）与 c））。e) 证明若 K 满足 d) 的条件，则 K 上每个代数元在 K 上的次数都形如 \(2^q\) （对某个 q）（仿照 VI 第 26 页命题 8，对所考虑元素的次数中 2 的指数作归纳）。另一方面，证明 K 的任何二次扩张都不是 A-可序的。由此推出 K 是极大 A-可序域（利用伽罗瓦理论和 I 第 77 页命题 12）。
\item
  设 K 是 A-可序域，且 \(\Omega\) 是 K 的代数闭扩张。证明存在一个包含于 \(\Omega\) 且包含 K 的极大 A-可序域。
\end{enumerate}

\begin{enumerate}
\def\labelenumi{\arabic{enumi})}
\setcounter{enumi}{8}
\tightlist
\item
  设 \(q > 0\) 为一个非平方自然数；令 \(A\) 表示集合 \(\{-1, \sqrt{q}\}\) （在域 \(K = Q(\sqrt{q})\) 中）；证明 \(K\) 是 \(A\) -可序的。证明存在一个扩张 \(E\) （ \(K\) 的），使得 \(E\) 是毕达哥拉斯的、可序的，并且 \(E\) 上每个奇数次多项式在 \(E\) 中有一个根，但 \(E\) 不具有极大序域的结构（考虑 \(A\) -可序的 \(K\) 的一个极大扩张）.
\end{enumerate}

¶ 10) a) 设 K 是一个可序域，E 是 K 的一个伽罗瓦扩张。证明：要么 E 是可序的，要么存在 K 的一个包含于 E 的可序代数扩张 F，使得 E 是 F 的二次扩张（使用 V 第 70 页定理 6）。

\begin{enumerate}
\def\labelenumi{\alph{enumi})}
\setcounter{enumi}{1}
\tightlist
\item
  证明多项式 \(X^{4} + 2\) 在 Q 上不可约，并且在 Q 上添加该多项式的一个根所得到的 Q 的 4 次扩张 K 不包含除 Q 以外的任何可序子域（利用伽罗瓦理论找出 K 的所有子域）。
\end{enumerate}

\begin{enumerate}
\def\labelenumi{\arabic{enumi})}
\setcounter{enumi}{10}
\tightlist
\item
  设 K 是一个序域，G 是 K 的一个子域。a) 证明：\(x \neq 0\) 关于 G 是无穷大的，当且仅当 \(x^{-1}\) 关于 G 是无穷小的。我们说 K 与 G 可比较，如果 K 中没有关于 G 无穷大的元素（从而也没有 K 的非零元素关于 G 是无穷小的）。K 与它的素子域 Q 可比较的充要条件是 K 是阿基米德的（VI，第～35 页，习题 31），* 从而 K 同构于 R 的一个子域（VI，第～36 页，习题 33c））。
\end{enumerate}

\begin{enumerate}
\def\labelenumi{\alph{enumi})}
\setcounter{enumi}{1}
\item
  证明在由 K 中关于 G 不是无穷大的元素组成的环 \(F(G)\) 中，关于 G 为无穷小的元素组成的集 \(I(G)\) 是一个极大理想；此外，商域 \(K(G) = F(G)/I(G)\) 上由 \(F(G)\) 的序诱导的序与环结构相容，并且是全序的。
\item
  证明模 \(I(G)\) 的一个等价类至多包含 G 的一个元素；由此推出从 G 到 \(K(G)\) 的自然映射是序域 G 到子域 \(G'\) （ \(K(G)\) 的）上的一个同构，并且 \(K(G)\) 与 \(G'\) 可比较。
\end{enumerate}

\begin{enumerate}
\def\labelenumi{\arabic{enumi})}
\setcounter{enumi}{11}
\item
  设 K 为一个极大序域，f 为 K 上的一个多项式，并设 a 和 b 为 f 的两个根，满足 a \textless{} b，且使 f 在 a 与 b 之间没有根。证明：若 g 是 K 上的有理函数，其分母在 \(a \leqslant x \leqslant b\) 上不消失，则方程 \(f(x)g(x) + f'(x) = 0\) 在 \([a, b]\) 中有奇数个解（其中每个解都按重数计算；使用 VI 第 25 页命题 5）。由此推出，如果 h 是 K 上的有理函数，以 a 和 b 为根，且其分母在 \([a, b]\) 中不消失，则方程 \(h'(x) = 0\) 在 \([a, b]\) 中至少有一个根。
\item
  设 K 是一个极大序域，h 是 K 上的一个有理函数，且 \([a, b]\) 是 h 在其上有定义的闭区间。证明存在 \(c \in ]a, b[\) 使得 \(h(b) - h(a) = (b - a)h'(c)\) （利用习题 12）。由此推出，h 在 \([a, b]\) 上是递增函数，当且仅当在此区间内 \(h'(x) \geqslant 0\) （要看出该条件是必要的，可用 \(h'(x) = 0\) 的根分割该区间）.
\item
  \begin{enumerate}
  \def\labelenumii{\alph{enumii})}
  \tightlist
  \item
    设 K 为一个极大序域，E 为 K 的子域，f 为 E 上的多项式；证明 f 在 K 中的所有根都属于 F(E)（习题 11，b））（利用 VI 第 24 页命题 4）。由此推出，若 G 为 K 的子域且 E ⊂ K 是与 G 可比较的扩张（习题 11），则 K 中在 E 上代数的元素之集是一个与 G 可比较的极大序域。
  \end{enumerate}
\end{enumerate}

\begin{enumerate}
\def\labelenumi{\alph{enumi})}
\setcounter{enumi}{1}
\item
  由a)推出，在a)的假设下，域K(G)（习题11，b)）是一个极大序域。
\item
  设 \(f\) 是次数 \(\geqslant 1\) 的 \(G\) 上的多项式。证明 \(f(t)\) 关于 \(G\) 是无穷大的，当且仅当 \(t\) 关于 \(G\) 是无穷大的；元素 \(f(t)\) 关于 \(G\) 是无穷小的，当且仅当 \(t\) 模 \(I(G)\) 同余于 \(f\) 在 \(K\) 中的一个根（把 \(K\) 用 \(f\) 与 \(f'\) 在 \(K\) 中的根分割成区间，并应用习题 13；注意若 \(x < t\) 且 \(x \in K\) 不模 \(I(G)\) 同余于 \(t\) ，则存在 \(y \in K\) 使得 \(x < y < t\) 且 \(y - x \in G\) ).
\end{enumerate}

\begin{enumerate}
\def\labelenumi{\arabic{enumi})}
\setcounter{enumi}{14}
\tightlist
\item
  设 E 是一个序域，K 是 E 的子域，使得 E 在 K 上是代数的，并且 R 是 K 的一个极大序扩张。
\end{enumerate}

\begin{enumerate}
\def\labelenumi{\alph{enumi})}
\item
  在 E - K 的元素中，设 \(x_0\) 是其在 K 上的极小多项式 \(f\) 的次数尽可能小的一个。证明 \(f'(T)\) 是 R{[}T{]} 中形如 \((T - a_i)^2 + c_i^2\) ( \(a_i, c_i \in \mathbb{R}\) ) 的因子与一次因子 \(T - b_j\) （其中 \(b_j\) 是 K 中两两不同的元素）的乘积。由此推出，存在 R 中的一个元素 \(y_0\) ，使得 \(f(y_0) = 0\) ，且 \(x_0 - z\) 与 \(y_0 - z\) 分别在 E 与 R 中对所有 \(z \in \mathbb{K}\) 具有相同的符号（参见 VI，第 43 页，习题 26，c））。
\item
  由 a) 推出，存在一个从序域 E 到 R 的某个子域上的 K-同构。（首先证明对 E 的子域 \(\mathrm{K}(x_{0})\) 存在这样的同构；对任意多项式 \(g \in \mathrm{K}[\mathrm{T}]\) （满足 \(\deg(g) < \deg(f)\) ），用 a) 中对 \(f'\) 用过的同样论证来求 \(g(x_{0})\) 与 \(g(y_{0})\) 的符号；然后应用佐恩引理）。
\item
  证明：若R和R'是K的两个极大序代数扩张，则存在恰好一个从R到R'的K-同构（应用b)）。
\end{enumerate}

\begin{enumerate}
\def\labelenumi{\arabic{enumi})}
\setcounter{enumi}{15}
\tightlist
\item
  设K是一个极大序域，并设
\end{enumerate}

\[
g (\mathrm{T}) = a _ {0} + a _ {1} \mathrm{T} + \dots + a _ {n} \mathrm{T} ^ {n}
\]

是 K{[}T{]} 中的一个非零多项式，其根全部位于 K 中。证明对于任意多项式 \(f \in K[T]\) ，多项式

\[
a _ {0} f (\mathrm{T}) + a _ {1} f ^ {\prime} (\mathrm{T}) + a _ {2} f ^ {\prime \prime} (\mathrm{T}) + \dots + a _ {n} f ^ {(n)} (\mathrm{T})
\]

的不属于 K 的根（按重数计）的个数至多等于 f 的不属于 K 的根（按重数计）的个数（对 n = 1 使用习题 12，然后归纳进行）。由此推出多项式

\[
a _ {0} + \frac {a _ {1}}{1 !} \mathrm{T} + \frac {a _ {2}}{2 !} \mathrm{T} ^ {2} + \dots + \frac {a _ {n}}{n !} \mathrm{T} ^ {n}
\]

的所有根都在K中。

\begin{enumerate}
\def\labelenumi{\arabic{enumi})}
\setcounter{enumi}{16}
\tightlist
\item
  设 K 是一个极大序域，g 是 K{[}T{]} 中的非零多项式，其根均在 K 中，且不属于 K 的区间 \([0, n]\) 。证明对于每一个多项式 \(f(T) = a_{0} + a_{1}T + \cdots + a_{n}T^{n}\) （属于 K{[}T{]}），多项式
\end{enumerate}

\[
a _ {0} g (0) + a _ {1} g (1) \mathrm{T} + a _ {2} g (2) \mathrm{T} ^ {2} + \dots a _ {n} g (n) \mathrm{T} ^ {n}
\]

的不属于 K 的根（按重数计）的个数至多等于 f 的不属于 K 的根（按重数计）的个数（方法与习题 16 相同）。

\begin{enumerate}
\def\labelenumi{\arabic{enumi})}
\setcounter{enumi}{17}
\item
  设 K 为一个极大序域，f 为 K{[}T{]} 中次数为 n 的多项式，其所有根都在 K 中。证明对 K 中所有 \(c \neq 0\) ，多项式 \(f^{2} + cf'\) 在 K 中按重数计至少有 n - 1 个且至多有 \(n + 1\) 个根（使用习题 13）。
\item
  设 K 是一个极大序域，\(f\) 是 K{[}T{]} 中的一个非零多项式。使得对每个多项式 \(g \in \mathbf{K}[\mathbf{T}]\) ，\(fg + g'\) 的不属于 K 的根（按重数计）的个数至多等于 \(g\) 的不属于 K 的根（按重数计）的个数，其必要且充分条件是 \(f(\mathbf{T}) = a - b\mathbf{T}\) ，其中 \(b \geqslant 0\) （在 K 中）。（要看出该条件是充分的，使用习题 12；要证明它是必要的，取 g = 1 和 g = f）。
\item
  设 \(\mathbf{K}\) 是一个极大序域，且
\end{enumerate}

\[
f (\mathrm{T}) = a _ {0} + \binom {n} {1} a _ {1} \mathrm{T} + \binom {n} {2} a _ {2} \mathrm{T} ^ {2} + \dots + a _ {n} \mathrm{T} ^ {n}
\]

是 K{[}T{]} 中的一个多项式。对于所有满足 \(0 \leqslant p < p + q \leqslant n\) 的整数 p, q，多项式

\[
a _ {p} + \binom{q}{1} a _ {p + 1} \mathrm{T} + \binom{q}{2} a _ {p + 2} \mathrm{T} ^ {2} + \dots + a _ {p + q} \mathrm{T} ^ {q}
\]

（按重数计）不属于 K 的根的个数至多等于 f 的（按重数计）不属于 K 的根的个数（利用习题 12）。

由此推出，如果 \(b_{1}, \ldots, b_{n}\) 是 K 中 n 个两两不同的 \textgreater0 的元素，且若

\[
\left(\mathrm{T} + b _ {1}\right) \dots \left(\mathrm{T} + b _ {n}\right) = \mathrm{T} ^ {n} + \binom {n} {1} m _ {1} \mathrm{T} ^ {n - 1} + \dots + m _ {n}
\]

则

\[
m _ {k} ^ {1 / k} > m _ {k + 1} ^ {1 / (k + 1)} \quad \text { for } \quad 1 \leqslant k \leqslant n - 1.
\]

\begin{enumerate}
\def\labelenumi{\arabic{enumi})}
\setcounter{enumi}{20}
\tightlist
\item
  设 \((a_{i})_{1 \leqslant i \leqslant n}\) 为序域 K 中 n 个元素组成的有限序列，且不全为零；设 \((a_{i_{k}})_{1 \leqslant k \leqslant p}\) 为 \((a_{i})\) 的由非零 \(a_{i}\) 组成的子序列 ( \(i_{1} < i_{2} < \cdots < i_{p}\) ）；则指标 \(k \leqslant p - 1\) （使得 \(a_{i_{k}}\) 与 \(a_{i_{k+1}}\) 具有相反符号者）的个数称为序列 \((a_{i})\) 的变号数。
\end{enumerate}

设 K 为一个极大序域，f 为 K{[}T{]} 中次数为 n \textgreater{} 0 的非零多项式；对所有 \(a \in K\) ，令 \(w(a)\) 表示序列 \((f^{(i)}(a))_{0 \leq i \leq n}\) 的变号数。若 v 是 f 在区间 \([a, b]\) (a \textless{} b) 内按重数计的根的个数，证明 \(v \leq w(a) - w(b)\) 且 \(w(a) - w(b) - v\) 是偶数（“布当–傅里叶法则”；用 f 的根分割区间 \([a, b]\) ，并计算当 x 经过其中一个根时 \(w(x)\) 的变号数的变化）。

由此推出，如果 \(f(T) = a_{0} + a_{1}T + \cdots + a_{n}T^{n}\) ，则 f 的按重数计的 \textgreater{} 0 的根的个数至多等于序列 \((a_{i})_{0 \leq i \leq n}\) 的变号数，并且这两个数之差是偶数（“笛卡儿法则”）。

\begin{enumerate}
\def\labelenumi{\arabic{enumi})}
\setcounter{enumi}{21}
\tightlist
\item
  设K是一个极大序域，并设
\end{enumerate}

\[
f (\mathrm{T}) = a _ {0} + a _ {1} \mathrm{T} + \dots + a _ {n} \mathrm{T} ^ {n}
\]

是 K{[}T{]} 中的多项式，满足 \(a_0 \neq 0 \neq a_n\) ，且

\[
a _ {p} = a _ {p + 1} = \dots = a _ {p + 2 m - 1} = 0 \quad (1 \leqslant p <   p + 2 m - 1 <   n).
\]

证明 \(f\) 按重数计至多有 \(n - 2m\) 个根在 \(\mathbf{K}\) 中（应用笛卡儿法则）。

由此推出，如果 \(g(T) = 1 + c_{1}T + \cdots + c_{n}T^{n}\) 的所有根都在 K 中，并且 \(h(T) = 1 + b_{1}T + \cdots + b_{2m}T^{2m}\) 是由前 \(2m + 1\) 项组成的多项式，这些项取自形式幂级数 \(1/g \in K[[T]]\) ，则 h 在 K 中没有根。

\begin{enumerate}
\def\labelenumi{\arabic{enumi})}
\setcounter{enumi}{22}
\tightlist
\item
  设 K 是一个极大序域，并设 \(a_{1}, \ldots, a_{n}, b_{1}, \ldots, b_{n}\) 是 K 的 2n 个元素，使得 \(\sum a_{i} \neq 0\) 且 \(b_{1} < b_{2} < \cdots < b_{n}\) 。证明对于任意整数 \(m \geqslant 1\) ，多项式在 K 中的根（按重数计）的个数 v
\end{enumerate}

\[
f (\mathrm{T}) = a _ {1} \left(\mathrm{T} - b _ {1}\right) ^ {m} + a _ {2} \left(\mathrm{T} - b _ {2}\right) ^ {m} + \dots + a _ {n} \left(\mathrm{T} - b _ {n}\right) ^ {m}
\]

至多等于序列的变号数w，

\[
\left(a _ {1}, a _ {2}, \dots , a _ {n}, (- 1) ^ {m} a _ {1}\right),
\]

且差w - v为偶数。（通过对n归纳证明，将习题13（VI，第40页）应用于形如 \(f(\mathrm{T})/(\mathrm{T}-c)^{m}\) .)

\begin{enumerate}
\def\labelenumi{\arabic{enumi})}
\setcounter{enumi}{23}
\item
  设 K 为一个极大序域；考虑有理函数域 K(X) 上的一个序，使其成为 K 的序扩张。证明该序由集合 A 决定，其中 A 为 \(x \in K\) 使得 x \textless{} X 的集合（利用 VI 第 28 页命题 9 证明 K 上每个多项式 f 的符号由 A 决定）。反之，证明对每个满足 \(y \in A\) 当 \(y \leqslant x\) 中，且 \(x \in A\) 时的集合 A ⊂ K，都存在 K(X) 的一个序，使其成为 K 的序扩张，并且 A 恰为满足 \(x \in K\) 使得 x \textless{} X 的集合（用同样方法）。当 A 有最大元，或 CA 有最小元，或 A = K 或 A = ∅ 时，对应的 K(X) 的序使得 K(X) 与 K 不可比较。另一方面，若 A 与 CA 均非空，且 A 无最大元、CA 无最小元，则 K(X) 与 K 可比较。
\item
  利用习题 24 找出一个域 E 的例子，使其具有多个序域结构，且其中任何一个都不能通过 E 的自同构由另一个诱导得到。由此找出一个域的例子，它带有与其环结构相容的序，但不是格序（考虑 E × E 中的对角线，并利用 VI 第 38 页习题 5）。
\item
  \begin{enumerate}
  \def\labelenumii{\alph{enumii})}
  \tightlist
  \item
    若 K 是阿基米德序域（VI 第 40 页习题 11 a)），且 x 与 y 是 K 中满足 x \textless{} y，证明存在 \(r \in Q\) 使得 x \textless{} r \textless{} y. * 推导出存在唯一一个与 K 同构的 R 的序子域。 *
  \end{enumerate}
\end{enumerate}

\begin{enumerate}
\def\labelenumi{\alph{enumi})}
\setcounter{enumi}{1}
\tightlist
\item
  考虑域 \(\mathbf{K} = \mathbf{Q}(\mathbf{X})\) 的序，其中 \(\mathbf{X}\) 关于 \(\mathbf{Q}\) 是无穷大的（习题 24）。证明 \(\mathbf{K}\) 与子域 \(\mathbf{Q}(\mathbf{X}^2)\) 可比较，并给出两个元素 \(x, y\) （属于 \(\mathbf{K}\) ）的例子，使得 \(x < y\) 且 \(\mathbf{Q}(\mathbf{X}^2)\) 中没有元素位于 \(x\) 与 \(y\) 之间。c) 设 \(\mathbf{K}\) 为在 \(b\) 中定义的序域；证明多项式
\end{enumerate}

\[
f (\mathbf {Y}) = (\mathbf {Y} ^ {2} - \mathbf {X}) (\mathbf {Y} ^ {2} - 4 \mathbf {X}) - 1
\]

在 K 上不可约，它在 K 的每个极大序扩张中都有根，并且 \(f(a) > 0\) 对所有 \(a \in K\) .

\begin{enumerate}
\def\labelenumi{\arabic{enumi})}
\setcounter{enumi}{26}
\tightlist
\item
  设 K 是一个序域，设 E 是 K 的一个纯超越扩张，并设 \((\mathbf{X}_{i})_{i\in I}\) 是 E 的一个超越基（V，第 109 页）。
\end{enumerate}

\begin{enumerate}
\def\labelenumi{\alph{enumi})}
\item
  * 若 K 是阿基米德的，则 E 具有 K 的序扩张结构，且 E 与 K 可比较，当且仅当 I 的基数至多等于 R 在 K 上的超越基的基数（其中 K 被视为 R 的序子域，参见习题 26 a)）；此时，此类结构的集合与从 I 到 R 的、使 \(f(I)\) 在 K 上代数无关的单射映射所成的集合成一一对应。*
\item
  如果K不是阿基米德的，那么在E上总是存在（至少）一种结构，使得E成为K的序扩张，并且E与K是可比较的（当I只包含一个元素时，使用习题24并注意到Q在K中没有上确界；通过佐恩引理过渡到一般情形）。
\end{enumerate}

* 28) a) 设 K 是 R 的一个子域，并且设 \(\theta\) 是在 K 上代数的实数。证明 K(θ) 上作为 K 的序扩张的结构的个数等于 \(\theta\) 的实共轭元的个数（利用习题 14，a）（VI，第 40 页）和 26，a））。

\begin{enumerate}
\def\labelenumi{\alph{enumi})}
\setcounter{enumi}{1}
\tightlist
\item
  设 K 为 R 的一个子域，带有 n 个不同的作为 Q 的序扩张的结构，且均为阿基米德的；设 \(P_{i}\) ( \(1 \leqslant i \leqslant n\) ) 为这些序下的正元素集合。证明存在一个由 K 中 n 个元素组成的族 \(b_{i}\) ( \(1 \leqslant i \leqslant n\) )，使得 \(b_{i} \in P_{i}\) 对所有 i 成立，且 \(b_{i} \notin P_{k}\) 对所有 \(k \neq i\) 成立。（对于每一对不同的指标 i, k，设 \(a_{ik} \in P_{k}\) 使得 \(a_{ik} \notin P_{i}\) ；考虑元素
\end{enumerate}

\[
c _ {i} = \sum_ {k \neq i} \left(\frac {1 + a _ {i k}}{1 - a _ {i k}}\right) ^ {2 r}
\]

，其中 \(r\) 是一个足够大的整数。）

\begin{enumerate}
\def\labelenumi{\alph{enumi})}
\setcounter{enumi}{2}
\tightlist
\item
  设 \(\mathbf{K} = \mathbf{Q}(\theta)\) 为 \(\mathbf{Q}\) 的一个包含于 \(\mathbf{R}\) 的代数扩张，并设 \(n\) 为 \(\theta\) 的实共轭元的个数。设 \(\mathbf{C}\) 为 \(\mathbf{K}\) 中元素的平方和所构成的集合；证明 \(\mathbf{C}^* = \mathbf{C} \cap \mathbf{K}^*\) 是 \(\mathbf{K}^*\) 的一个子群，且 \((\mathbf{K}^*: \mathbf{C}^*) = 2^n\) 。（使用 \(b\) ）以及 VI 第 23 页推论 2。）*
\end{enumerate}

\begin{enumerate}
\def\labelenumi{\arabic{enumi})}
\setcounter{enumi}{28}
\item
  设 K 为一个极大序域，G 为 K 的子域。证明：包含于 K 中且与 G 可比较的 G 的扩张之集是归纳有序的；若 \(E_{0}\) 是该集合的一个极大元素，证明 \(E_{0}\) 同构于 VI 第 40 页习题 11, b) 中定义的域 \(\mathrm{K}(\mathrm{G})\) （证明从 \(\mathrm{F}(\mathrm{G})\) 到 \(\mathrm{K}(\mathrm{G})\) 上的自然映射把 \(E_{0}\) 映到 \(\mathrm{K}(\mathrm{G})\) 上：首先利用 VI 第 40 页习题 14, a) 证明 \(E_{0}\) 是一个极大序域，进而利用习题 24 证明 \(\mathrm{K}(\mathrm{G})\) 中没有任何元素在 \(E_{0}\) 的像上是超越的）。
\item
  \begin{enumerate}
  \def\labelenumii{\alph{enumii})}
  \tightlist
  \item
    设 K 是一个极大序域，并设 m 和 M 是 K 中满足 m \textless{} M 的两个元素。证明每个多项式 \(f \in K[X]\) （在区间 \((m, M)\) 内为正的）都是形如 \((\alpha X + \beta) g^{2}\) 的多项式之和，其中 \(g \in K[X]\) ，且 \(\alpha X + \beta\) 在 \((m, M)\) 内为正。（化归为一次或二次多项式的情形，并利用公式
  \end{enumerate}
\end{enumerate}

\[
\begin{array}{l} (\mathbf {X} - a) (\mathbf {X} - b) = (\mathbf {X} - b) ^ {2} + (b - a) (\mathbf {X} - b) \\ (\mathbf {X} - a) (b - \mathbf {X}) = ((\mathbf {X} - a) (b - \mathbf {X}) ^ {2} + (b - \mathbf {X}) (\mathbf {X} - a) ^ {2}) / (b - a) \end{array}
\]

对 \(a < b\) .)

\begin{enumerate}
\def\labelenumi{\alph{enumi})}
\setcounter{enumi}{1}
\tightlist
\item
  证明当 K 为任意序域时，a) 的结果不再必然成立（注意一个多项式在 K 中可以为 \textgreater0，但在 K 的一个极大序扩张中不为正；参见习题 26，c)）。
\end{enumerate}

\begin{enumerate}
\def\labelenumi{\arabic{enumi})}
\setcounter{enumi}{30}
\tightlist
\item
  \begin{enumerate}
  \def\labelenumii{\alph{enumii})}
  \tightlist
  \item
    设 \(\mathbf{K}\) 是一个域，其代数闭包 \(\mathbf{E}\) 是 \(\mathbf{K}\) 的素数次 \(q\) 扩张。证明 \(\mathbf{K}\) 是完全域（V，第 7 页）。
  \end{enumerate}
\end{enumerate}

\begin{enumerate}
\def\labelenumi{\alph{enumi})}
\setcounter{enumi}{1}
\item
  证明 q 不同于 K 的特征（V，第 162 页，习题 9）。
\item
  证明K包含q次单位根，且E是形如 \(X^{q}-a\) 的不可约多项式的分裂域；由此推出q=2（否则，利用习题5，V，第161页推出 \(X^{q^{2}}-a\) 将是不可约的）。同时证明 -a 在 K 中是平方（见 V，前引），-1 在 K 中不是平方，并且 \(E = K(i)\) ( \(i^{2} = -1\) ).
\item
  现在假设 K 的代数闭包 E 在 K 上具有任意有限次数 \(\neq1\) 。证明 \(i\notin K\) 且 \(\mathrm{E}=\mathrm{K}(i)\) （若 \(\mathrm{E}\neq\mathrm{K}(i)\) ，用伽罗瓦理论证明会存在一个域 F 使得 \(\mathrm{K}(i)\subset\mathrm{F}\subset\mathrm{E}\) ，并且 E 在 F 上的次数为素数；然后应用 c)）。由此推出 K（在某个适当的序下）是一个极大序域（用归纳法证明 K 中任意 n 个平方的和是 K 中的平方；为了证明 \(a^{2}+b^{2}\) 是平方，考虑平方根 \(x+iy\) （ \(a+ib\) 在 \(\mathrm{K}(i)\) 中的）).
\end{enumerate}

\begin{enumerate}
\def\labelenumi{\arabic{enumi})}
\setcounter{enumi}{31}
\tightlist
\item
  设 A 是一个特征为 0 的代数闭域。
\end{enumerate}

\begin{enumerate}
\def\labelenumi{\alph{enumi})}
\item
  在群 \(\operatorname{Aut}(\mathbf{A})\) （ \(\mathbf{A}\) 的自同构群）中，仅有的有限阶元素是单位元和二阶元素（ \(\mathbf{A}\) 的对合）；后者与极大序子域 \(\mathbf{E}\) （ \(\mathbf{A}\) 中满足 \([\mathbf{A}:\mathbf{E}] = 2\) 者）一一对应，即 \(\mathbf{A} = \mathbf{E}(i)\) （参见习题 31）。若 \(\sigma\) 是 \(\mathbf{A}\) 的一个对合，且 \(\mathbf{E}\) 是 \(\sigma\) 的固定子域，则群 \(\operatorname{Aut}(\mathbf{E})\) （ \(\mathbf{E}\) 的自同构群）同构于商 \(Z(\sigma) / \{e, \sigma\}\) ，其中 \(Z(\sigma)\) 是 \(\sigma\) 在 \(\operatorname{Aut}(\mathbf{A})\) 中的中心化子。
\item
  如果 A 在其素域 Q 上是代数的，那么 A 的任意两个对合在 Aut(A) 中共轭，且对合 \(\sigma\) 的中心化子是 \(\{e, \sigma\}\) ；群 Aut(A) 与 A 的对合集合具有连续统的基数，且 Aut(A) 的中心是 \(\{e\}\) .
\item
  若 A 在 Q 上具有有限的超越次数，证明中心化子 \(Z(\sigma)\) （其中 \(\sigma\) 是 A 的任意一个对合）是可数的（注意 \(\sigma\) 的固定子域 E 是可数的，并且 E 的任何自同构，只要固定 E 在 Q 上的某个超越基的元素，就是恒等映射）。如果 deg. \(tr_{Q}A = n\) ，证明至少存在 \(n + 1\) 个 A 的对合的共轭类（若 \(E \subset A\) 是极大可序的且 \([A : E] = 2\) ，考虑 E 中由 Aut(E) 固定的元素所构成的子域在 Q 上的超越次数，并利用习题 24）。
\item
  若 A 在 Q 上具有无限超越次数 m，证明存在 A 的对合 \(\sigma\) ，使得 \(\operatorname{Card}(Z(\sigma)) = \operatorname{Card}(\operatorname{Aut}(A)) = 2^{m}\) （利用习题 24 证明存在极大可序域 E，使得 deg. \(tr_{Q}E = m\) ，使得对于 E 在 Q 上的任意超越基 B 的任意子集 M，存在 E 的一个自同构 \(u_{M}\) 满足： \(u_{M}(x) \neq x\) （对 \(x \in M\) ），且 \(u_{M}(x) = x\) （对 \(x \in B - M\) ）).
\item
  如果 \(\deg \cdot tr_{Q}A = 1\) ，证明 A 的包含无限中心化子的对合只有一个共轭类。（若 E 是 Q 的超越次数为 1 的极大可序扩张，且 \(\operatorname{Aut}(E)\) 非平凡，则利用习题 24 证明 E 不能与 Q 可比；若 \(t \in E\) 关于 Q 是无限大的，则证明对 E 的任何自同构 u， \(u(t)\) 也是无限大的，而反之，若 \(t'\) 关于 Q 是无限大的，则存在 E 的自同构 u 使得 \(u(t) = t'\) ，利用 VI 第 40 页习题 15。）
\item
  设 u 是 A 的一个自同构，且 u 与 A 的所有对合交换。证明若 \(a \in A\) 在 Q 上是超越的，则 a 和 \(u(a)\) 在 Q 上代数相关。（考虑 A 中两个元素 b, c，使得 \(b^{2} = a\) , \(c^{2} = -u(a)\) 并观察到存在一个极大的可序扩张 \(E \subset A\) 使得 \([A : E] = 2\) 且 b, c 属于 E；注意到 \(u(E) = E\) 便得出矛盾。）
\end{enumerate}

\(g\) ）由 \(f\) ) 推出 \(u(a) = a\) 对每个元素 \(a \in \mathbf{A}\) （它在 \(\mathbf{Q}\) 上是超越的）成立（考虑 \(\mathbf{Q}(a)\) 的一个极大序扩张 E，它包含于 A 中，在 \(\mathbf{Q}(a)\) 上是代数的，且与 \(\mathbf{Q}\) 可比，利用习题 24）。由此得出结论： \(u\) 必然是恒等映射（注意若 \(a\) 在 \(\mathbf{Q}\) 上是代数的，且 \(b\) 在 \(\mathbf{Q}\) 上是超越的，则 \(ab\) 在 \(\mathbf{Q}\) 上是超越的）.

\begin{enumerate}
\def\labelenumi{\arabic{enumi})}
\setcounter{enumi}{32}
\tightlist
\item
  设 \(\mathbf{K}\) 是一个交换域（特征任意， \(q\) )，并令 \(l\) 是一个素数。设 \(\mathbf{K}'\) 是 \(\mathbf{K}\) 的一个代数扩张，满足以下两个性质：
\end{enumerate}

\begin{enumerate}
\def\labelenumi{(\roman{enumi})}
\item
  \(\mathbf{K}[\mathbf{X}]\) 中每个次数不被 \(l\) 整除的非常数多项式在 \(\mathbf{K}'\) 中有一个根 ;
\item
  \(\mathbf{K}'[\mathbf{X}]\) 中每个次数等于 \(l\) 的多项式在 \(\mathbf{K}'\) 中有一个根 .
\end{enumerate}

\begin{enumerate}
\def\labelenumi{\alph{enumi})}
\item
  证明完全闭包 \(\mathbf{K}_1\) （ \(\mathbf{K}\) 的）包含于 \(\mathbf{K}'\) ，并且满足与 (i) 类似的条件。
\item
  K 的每一个次数不被 l 整除的代数扩张都同构于 \(K'\) 的某个子扩张（使用本原元素定理）。
\item
  设 L 是 \(K'\) 的一个包含 K 的子域，并设 M 为 L 的伽罗瓦扩张，其次数为 l 的幂。证明 M 同构于 \(K'\) 的一个子扩张（利用 l-群是幂零群这一事实）。
\end{enumerate}

\(d\) ）由此推出 \(\mathbf{K}'\) 是 \(\mathbf{K}\) 的一个代数闭包（若 \(\mathbf{F}\) 是 \(\mathbf{K}\) 的一个有限次数的（可分）扩张，考虑一个有限次数的伽罗瓦扩张 \(\mathrm{F}_1\supset \mathrm{F}\) ，以及 \(\mathrm{F}_1\) 在西罗 \(l\) -子群（ \(\mathrm{F}_1\) 在 \(\mathbf{K}\) 上的伽罗瓦群的）作用下的固定子域）.

\begin{enumerate}
\def\labelenumi{\arabic{enumi})}
\setcounter{enumi}{33}
\tightlist
\item
  设 E 是一个全序集。一个（戴德金）分割是指 E 的一个划分 (S, D)，其中 S 和 D 是两个非空子集，且 D 具有形式 M(A)，即 E 的某个非空子集 A 的上界集；D 有最小元 x 的充分必要条件是 \(D = \{x, \rightarrow\{, \text{所以 } S = \} \leftarrow, x\{.\)
\end{enumerate}

在全序阿贝尔群 G 中，一个分割 (S, D) 称为真分割，如果对所有 e \textgreater{} 0（在 G 中）存在 \(x' \in S\) 与 \(x'' \in D\) 使得 \(x'' - x' < e\) 。证明 (S, D) 是真分割当且仅当 D 在由 G 的上界子集组成的幺半群 \(\mathfrak{M}(G)\) （VI，第 35 页，习题 30）中可逆；所有分割都是真分割当且仅当 G 是阿基米德的（VI，第 35 页，习题 31）；因此 * 同构于 R 的一个子群（VI，第 36 页，习题 33）。 *

\begin{enumerate}
\def\labelenumi{\arabic{enumi})}
\setcounter{enumi}{34}
\tightlist
\item
  序域 E 的子域 K 在 E 中稠密，如果对于每个非空开区间 \(a, b \in E\) 存在 \(x \in K\) 使得 a \textless{} x \textless{} b. 证明此时 E 与 K 是可比较的（习题 11，a））。若取 E 为 R(X)，其中 R 是极大序域，且 X 相对于 R 是无穷大的（VI，第 43 页，习题 24），证明 E 与子域 \(K = R(X^{2})\) 是可比较的，但 K 在 E 中并不稠密。
\end{enumerate}

要使 K 在 E 中稠密，其充要条件是：对所有 \(x \in E\) ，由 \(K \cap J \leftarrow, x\) 与 \(K \cap [x, \rightarrow]\) 组成的 K 的分割是一个真分割（习题 34）。

\begin{enumerate}
\def\labelenumi{\arabic{enumi})}
\setcounter{enumi}{35}
\item
  设 K 是一个序域，并设 \(\tilde{K}\) 为所有使得分割 (S, D) 为真分割的上界集 D 所构成的集合；证明 \(\tilde{K}\) 是一个序域，其中加法取为幺半群 \(\mathfrak{M}(K)\) 的加法，而两个元素 \(D_{1}\) 与 \(D_{2}\) （均属于 \(\tilde{K}\) 且包含于 \(K_{+}\) ）的乘积定义为集合 \(\langle D_{1}D_{2}\rangle\) （VI，第 35 页，习题 30），其中 \(D_{1}D_{2}\) 表示由 \(x_{1}x_{2}\) （ \(x_{1}\in D_{1}\) ， \(x_{2}\in D_{2}\) ）组成的集合。证明：对于从 K 到序域 E 的稠密子域 F 上的任意递增同态 f，存在从 E 到 \(\tilde{K}\) 中的递增同态 g，使得 \(g\circ f\) 是 K 到 \(\tilde{K}\) 中的自然单射。
\item
  设 K 是一个序域，且设 n \textgreater{} 0 为整数，使得 K 的每个元素 x \textgreater{} 0 都等于一个幂 \(y^{n}\) ，其中 y \textgreater{} 0 。证明域 \(\tilde{K}\) 具有相同的性质。
\item
  假设域 K 是极大序的。证明该域 \(\tilde{\mathbf{K}}\) 也是极大序的。（可以通过反证法来论证：考虑 \(\tilde{\mathbf{K}}\) 的一个极大序的序代数扩张 E，并假设 \(\mathrm{E} \neq \tilde{\mathbf{K}}\) 。那么，在 \(\mathrm{E} \cap [\tilde{\mathbf{K}}\) 的元素中存在一个元素 \(w\) ，其极小多项式 \(f \in \tilde{\mathbf{K}}[\mathbf{X}]\) 的次数最小，为 \(n > 1\) . 证明存在一个区间 \(\mathbf{a}, b\) （在 E 中），使得 \(a, b \in \mathbf{K}\) 且 \(a < w < b\) ，以及 E 的一个元素 \(e > 0\) ，使得 \(f'(x) > e\) 对所有 \(x \in [\mathbf{a}, b]\) 成立，或 \(f'(x) < -e\) 对所有 \(x \in [\mathbf{a}, b]\) 成立。证明存在 K 中的 \(r > 0\) ，使得 \(f(a) \leqslant -r\) 且 \(r \leqslant f(b)\) ，或 \(f(a) \geqslant r\) 且 \(f(b) \leqslant -r\) 。用 K 的元素逼近 \(f\) 的系数，并证明这样得到一个多项式 \(g \in \mathbf{K}[\mathbf{X}]\) ，使得函数 \(x \mapsto g(x)\) 在 \([\mathbf{a}, \mathbf{b}]\) 上单调，并且在此区间内 K 的某一点处取零值，该点任意接近 \(w\) .)
\item
  设 K 是一个序域，E 是一个极大序的序代数扩张。证明域 E 的 K-自同构只有恒等映射。（首先证明从 E 到自身的每个 K-同态都是满射，并且是序域的同构；如果 \(x \in E\) 具有极小多项式 \(f \in K[X]\) ，那么 E 的自同构必然置换属于 E 的 f 的根。由此推出它固定 x。）
\item
  设 K 为序域。把形式幂级数域 \(E = K((X))\) 构造成一个序域，取元素 \textgreater0 为这样的形式幂级数 \(\sum_{n=-h}^{\infty} r_n X^n\) ：其
\end{enumerate}

最低次数的非零系数为 \textgreater0.

\begin{enumerate}
\def\labelenumi{\alph{enumi})}
\item
  证明如果 \(L = K(X)\) 是 E 的子域，由单未定元的有理函数组成，则 \(E = \tilde{L}\) （习题 36）在序域的同构意义下（注意对 L 中每个元素 a \textgreater{} 0，存在整数 n \textgreater{} 0 使得 \(0 < X^{n} < a\) ).
\item
  假设 K 是极大序的。证明诸域 \(\mathbf{K}((\mathbf{X}^{1/n}))\) 的并集 F（按与 E 相同的方式排序）是 E 的一个序代数扩张，并且是极大序的（利用 V 第 150 页习题 2）。
\item
  证明域 \(\widetilde{\mathbf{F}}\) （习题 36）是由形如 \(\sum_{r} \alpha_r X^r\) 的形式级数组成的域，其中 \(r\) 遍历
\end{enumerate}

有理数集 Q，其中 \(\alpha_{r} \in K\) ，且使得 \(\alpha_{r} \neq 0\) 的 r 所成的集合是一个递增序列，该序列要么有限，要么趋于 \(+\infty\) ；按照与 E 相同的方式对该域进行排序。特别地 \(\widetilde{F} \neq F\) .

\begin{enumerate}
\def\labelenumi{\arabic{enumi})}
\setcounter{enumi}{40}
\tightlist
\item
  设 K 为一个序域，E 为 K 上维数为 2 的向量空间。
\end{enumerate}

\begin{enumerate}
\def\labelenumi{\alph{enumi})}
\item
  证明两两不同的直线三元组（ \(D_{1}\) , \(D_{2}\) , \(D_{3}\) ，均经过 0）的集合在群 \(GL^{+}(E)\) 的作用下有两条轨道（考虑固定两条直线的子群）。
\item
  两两不同的半直线三元组（ \(\Delta_{1}\) , \(\Delta_{2}\) , \(\Delta_{3}\) ，以 0 为原点）的集合在 \(GL(E)\) 的作用下有 7 条轨道，而在 \(GL^{+}(E)\) 的作用下有 14 条轨道（同样的方法）。
\end{enumerate}

\chapter{主理想整环上的模}

\section{主理想整环}

\subsection{主理想整环的定义}

回顾（I，第 104 页），交换环 A 的一个理想称为主理想，如果它具有形式 \((a) = A a\) ，对某个 \(a \in A\) .

定义1. --- 主理想整环是一个整环（I，第116页），其中每个理想都是主理想。

例子. --- 有理整数环 Z 是一个主理想整环（I，第 111 页）。若 K 是一个交换域，则 K 上一个未定元的多项式环 K{[}X{]} 是主理想整环（IV，第 11 页，命题 11）；形式幂级数环 K{[}{[}X{]}{]} 同样如此，因为该环的每个理想都具有形式 (X \(^{n}\) ) （IV，第 38 页，命题 12）。* p 进域中的整数环是一个主理想整环。*

如果 \(\mathbf{Q}(i)\) 表示由有理数域 Q 添加不可约多项式 \(X^{2}+1\) 的一个根 i 所得到的域，则形如 \(a+bi\) 的元素（ \(\mathbf{Q}(i)\) 中的，其中 a 和 b 是有理整数）构成 \(\mathbf{Q}(i)\) 的一个子环 A，称为「高斯整数」，它是一个主理想整环（VII，第 50 页，习题 7）。相比之下，在域 \(\mathbf{Q}(\rho)\) （其中 \(\rho\) 是 \(X^{2}+5\) 的一个根）中，由元素 \(a+b\rho\) （其中 a 和 b 为有理整数）构成的子环 B 不是主理想整环（VII，第 51 页，习题 12）。

域 K 上的两个未定元的多项式环 K{[}X, Y{]} 不是主理想整环。事实上，非零常数是仅有的同时整除 X 和 Y 的元素，而它们中没有一个能生成由 X 和 Y 生成的理想。

\subsection{主理想整环中的整除性}

设 A 为主理想整环，并设 K 为其分式域（I，第 116 页）；我们将看到由 K 的主分式理想（VI，第 6 页）组成的序群 \(P^{*}\) 是格序的；更精确地说：

命题1. --- 设 K 为主理想整环 A 的分式域，且设 \((x_{\iota})_{\iota \in I}\) 是 K 中具有公分母 \(b \in K^{*}\) 的元素族（换句话说， \(bx_{\iota} \in A\) 对所有 \(\iota\) ). 那么：

\begin{enumerate}
\def\labelenumi{\alph{enumi})}
\item
  族 \((x_{\iota})\) 在 K 中有一个最大公因子。
\item
  \((x_{i})\) 的每个最大公因子都可以表示为 \(d = \sum_{i}a_{i}x_{i}\) ，其中
\end{enumerate}

\(a_{i}\) 是A的元素，其中除有限多个外均为零。

事实上，理想 \(\sum_{i} A b x_{i}\) （ \(A\) 的）是主理想，因此具有形式 \(A d'\) 。令

\(d' = bd (d \in \mathbf{K})\) 。由关系式 \(d' = \sum_{i} a_i b x_i\) ，我们推导出 \(d = \sum_{i} a_i x_i\) ，其中

\(a_{i} \in A\) 。因此， \(x_{i}\) 的每个公因子都整除 d。另一方面，由于按构造 bd 是 \(bx_{i}\) 的公因子，由此可知 d 是 \(x_{i}\) 的公因子.

注：--- 命题 1 无限制地适用于 A 的元素的任意族 \((x_{i})\) （取 \(b = 1\) ），也适用于 K 中元素的任意有限族 \((x_{i})\) （若 \(x_{i} = c_{i}b_{i}^{-1}\) ，其中 \(c_{i} \in \mathbf{A}\) ， \(b_{i} \in \mathbf{A}\) ，则取 \(b\) 为诸 \(b_{i}\) 的乘积）.

推论。——设 \((x_{t})\) 是主理想整环 A 中元素的一个任意族，其中 A 作为子环包含在整环 B 中，并设 d 为族 \((x_{t})\) 在 A 中的一个最大公因子。则族 \((x_{t})\) 在 B 中有最大公因子，且 d 是其中之一。

确实， \(d\) 是 \(x_{i}\) 在 B 中的公因子。另一方面，关系式 \(d = \sum_{i} a_{i} x_{i}\) 表明 \(x_{i}\) 在 B 中的每个公因子都整除 \(d\) .

这个推论的一个重要应用是当 \(\mathbf{A} = \mathbf{K}[\mathbf{X}]\) 且 \(\mathbf{B} = \mathbf{E}[\mathbf{X}]\) 时，其中 \(\mathbf{K}\) 是一个域，且 \(\mathbf{E}\) 是 \(\mathbf{K}\) 的一个扩张（IV，第 12 页，推论 1）。

命题1的第一个断言表明序群 \(P^{*}\) 是格序的（VI，第 10 页）。特别地，K 的元素的每个有限族都有最小公倍元。因此我们可以将VI第10至17页中记为（DIV）的结果应用于主理想整环。

以下结果是命题1第二个断言的推论：

定理1（“贝祖恒等式”）。--- 主理想整环 A 中的元素 \(x_{\iota}\) ( \(\iota \in I\) ）是整体互素的，当且仅当存在 A 的元素 \(a_{\iota}\) ( \(\iota \in I\) ）（除有限多个外均为零），使得 \(\sum a_{\iota}x_{\iota} = 1\) .

该条件的必要性正是命题 1。反之，如果 \(\sum_{i} a_i x_i = 1\) ，那么 \(x_i\) 的每个公因子都整除 1，因此 1 是 \(x_i\) 的一个最大公因子.

命题 2. --- 设 \(a, b, d, m\) 与 \(p\) 是分式域 \(K\) （主理想整环 \(A\) 的）的元素.

\begin{enumerate}
\def\labelenumi{\alph{enumi})}
\item
  “d 是 a 和 b 的一个最大公因子”等价于“(d) = (a) + (b)”。
\item
  “m 是 a 和 b 的一个最小公倍元”等价于“m = (a) ∩ (b)”。
\item
  «p是A的不可约元素»等价于«(p)是A的非零极大理想»，也等价于«(p)是A的非零素理想»。
\end{enumerate}

我们已经证明了 \(a\) )（命题 1）。由于 \(a\) 与 \(b\) 的公倍数就是 \((a) \cap (b)\) 的元素，并且由于根据假设，理想 \((a) \cap (b)\) 是主理想，设 \((a) \cap (b) = (m)\) ，由此可得 \(m\) 是 \(a\) 与 \(b\) 的一个最小公倍元，这证明了 \(b\) )。最后，说 \(p \neq 0\) 是 \(A\) 的不可约元素，按定义（VI，第 17 页）意味着 \((p)\) 是 \(\neq A\) 的主理想（ \(A\) 中的）按包含关系排序所成族的一个极大元；由于 \(A\) 的理想都是主理想，这意味着 \((p)\) 是 \(A\) 的一个极大理想，从而得到 \(c\) )，根据 VI 第 17 页的注记。

在主理想整环 A 中，有限个理想的加法（相应地，交）也称为这些理想的 gcd（相应地，lcm）。

命题 3. --- 设 \(a, b, c\) 是主理想整环 A 的分式域的元素，并令 \(d\) 为 \(a\) 与 \(c\) 的一个最大公因子；则同余式 \(ax \equiv b \pmod{c}\) 有解 \(x_0 \in A\) 当且仅当 \(d\) 整除 \(b\) ；在这种情况下，A 的元素 \(x\) （满足 \(ax \equiv b \pmod{c}\) 者）恰好是那些满足 \(x \equiv x_0 \pmod{cd^{-1}}\) 的元素.

如果 \(ax \equiv b \pmod{c}\) ，其中 \(x \in \mathbf{A}\) ，则存在 \(y \in \mathbf{A}\) 使得 \(b = ax + cy\) ，所以 \(d\) 整除 \(b\) 。反之，若 \(d\) 整除 \(b\) ，则 \(b = ax_0 + cy_0\) ，其中 \(x_0, y_0 \in \mathbf{A}\) （命题 1），因此 \(ax_0 \equiv b \pmod{c}\) ；此外，关系 \(ax \equiv b \pmod{c}\) 便等价于 \(a(x - x_0) \equiv 0 \pmod{c}\) ；令 \(a = da'\) ， \(c = dc'\) ，后者变为 \(a'(x - x_0) \equiv 0 \pmod{c'}\) 。但这（对于 \(x \in \mathbf{A}\) ）等价于 \(x - x_0 \equiv 0 \pmod{c'}\) ，因为 \(a'\) 与 \(c'\) 互素（VI，第14页，命题10（DIV）及VI，第15页，命题11的推论2（DIV））。

命题4. --- 设 \((a_i)_{1 \leq i \leq n}\) 是主理想整环 A 中两两互素的元素组成的有限族。则从 \(\mathbf{A}/\left(\prod_{i=1}^{n} a_i\right)\) 到乘积 \(\prod_{i=1}^{n} \mathbf{A}/(a_i)\) 的典范同态（I，第 110 页）是 A-代数的同构。

这由I，第110页，命题9以及命题2 a)得出。

命题4的结论在未假设 \(a_{i}\) 两两互素的情况下不成立（参见VII，第24页，命题9）。

\subsection{在主理想整环中分解为不可约因子}

我们现在将 VI 第 18 页中关于分解为不可约元素的结果应用于主理想整环。根据命题 2，A 的一个元素 \(p \neq 0\) 是不可约的，当且仅当环 \(\mathrm{A}/(p)\) 是一个域（I，第 115 页，推论 1），即当且仅当同余式 \(ax \equiv b \pmod{p}\) 对每个 \(b \in \mathrm{A}\) 以及每个 \(a \in \mathrm{A}\) （不是 \(p\) 的倍元）在 A 中有解.

定义2. --- 设 A 是一个整环。A 的不可约元素的一个代表系是 A 的不可约元素的一个族 \((p_{\alpha})\) ，使得 A 的每个不可约元素恰好是其中一个 \(p_{\alpha}\) 的相伴元.

定理2. --- 设 A 是一个主理想整环，且设 \((p_{\alpha})\) 为 A 的不可约元素的一个代表系。则A的分式域中每个非零元素x都可以唯一地表示为如下形式

\[
x = u \prod_ {\alpha} p _ {\alpha} ^ {n _ {\alpha}},\tag{1}
\]

其中 u 是 A 的一个可逆元，且诸 \(n_{\alpha}\) 都是整数，除有限多个外均为零。要使 x 属于 A，必要且充分的条件是所有 \(n_{\alpha}\) 都是正的。

我们将使用关于分解为不可约元素之和的定理（VI，第 18 页，定理 2），上述陈述只是该定理的一种翻译。由于 \(\mathcal{P}^*\) 是一个格序群，为了验证该定理的假设确实成立，我们只需证明A的主理想的每个非空集合都包含一个极大元；而这由下面的引理得出：

引理1. --- 设 A 是一个环，使得 A 的每个左理想都是有限生成的。那么 A 的左理想组成的、按包含关系排序的每个非空集合 \(\Phi\) 都具有一个极大元。

由佐恩引理（集合论，III，第 154 页，定理 2）可知，只需证明 \(\Phi\) 是归纳的。现在如果 \((\mathfrak{a}_{\lambda})\) 是 \(\Phi\) 中元素的一个全序族，那么并集 \(\mathfrak{a}\) （诸理想 \(\mathfrak{a}_{\lambda}\) 的）是 \(\mathbf{A}\) 的左理想，因此具有有限生成系 \((a_i)_{1 \leqslant i \leqslant n}\) 。由于每个 \(a_i\) 属于某个理想 \(\mathfrak{a}_{\lambda_i}\) ，并且由于族 \((\mathfrak{a}_{\lambda})\) 是全序的，诸 \(a_i\) 都属于诸理想 \(\mathfrak{a}_{\lambda_i}\) 中最大的一个，设为 \(\mathfrak{a}_{\mu}\) 。于是 \(\mathfrak{a} = \mathfrak{a}_{\mu}\) 属于 \(\Phi\) ，后者因此确实是一个归纳集。

稍后我们将研究那些环 B，称为诺特环，使得 B 的每个非空理想集都包含一个极大元。

注记。 --- 族 \((u, (n_{\alpha}))\) 称为 x 分解为不可约因子的分解；滥用语言地，我们也称公式 (1) 是 x 分解为不可约因子的分解。如果 \(x = u \prod_{\alpha} p_{\alpha}^{n_{\alpha}}\) 与 \(y = v \prod_{\alpha} p_{\alpha}^{m_{\alpha}}\) 是

是 x 和 y 分解为不可约因子的分解，那么 x 整除 y 的一个必要且充分条件是 \(n_{\alpha} \leqslant m_{\alpha}\) 对所有 \(\alpha\) ；由此我们推导出以下公式

\[
\operatorname * {g c d} (x, y) = \prod_ {\alpha} p _ {\alpha} ^ {\inf (n _ {\alpha}, m _ {\alpha})}\tag{2}
\]

\[
\operatorname{lcm} (x, y) = \prod_ {\alpha} p _ {\alpha} ^ {\sup \left(n _ {\alpha}, m _ {\alpha}\right)}.\tag{3}
\]

定理2所表达的性质对于比主理想整环更广泛的一类环也成立；我们稍后将把它们作为唯一分解整环来研究；并且我们将看到，在任意多个未定元上的多项式环和形式幂级数环都是唯一分解整环（交换代数，第七章，第3节）。

\subsection{有理整数的整除性}

正如第 1 节所指出的，有理整数环 Z 是一个主理想整环；其分式域为 Q。Z 的可逆元乘法群 U 包含两个元素 1 和 -1。群 \(Q_{+}^{*}\) （由 \textgreater0 的有理数组成）恰好包含来自 Q 的每个相伴元素类中的一个元素；因此它与 Q 的主分式理想组成的乘法群 \(P^{*} = Q^{*}/U\) 同构，通常把它们等同看待。特别地，当在域 Q 中（相对于环 Z）使用 gcd 或 lcm 时，应理解为这些是 \(\geqslant 0\) 的元素；这一约定使我们能够谈论一族有理数的最大公因子和最小公倍元。

Z 中的不可约整数 \textgreater0 正是我们称之为素数的那些数（I，第 50 页）（它们有时被称为有理素数）；Z 的每个不可约元素因此具有 p 或 -p 的形式，其中 p 是素数，且素数集合 P 是 Z 的不可约元素的一个代表系。

命题5. --- 素数集合是无限的。

事实上，给定任意有限族 \((p_{i})\) \((1 \leqslant i \leqslant n)\) （由互不相同的素数组成），数 \(\left(\prod_{i=1}^{n} p_{i}\right) + 1\) （它是 \textgreater1 的）的任何素因子 q 与所有 \(p_{i}\) 不同，否则它将整除 1。

\subsection{一元多项式在域上的整除性}

交换域 K 上一个未定元的多项式环 K{[}X{]} 是主理想整环（IV，第 11 页，命题 11）。其分式域是系数取自 K 的、关于未定元 X 的有理函数域 K(X)。环 K{[}X{]} 包含次数为 0 的多项式构成的子环，即常数域，它与 K 等同；K* 中的元素在 K 中可逆，从而在 K{[}X{]} 中也可逆；反之，公式 \(\deg(uv) = \deg(u) + \deg(v)\) 表明每个可逆多项式都有次数 0；K{[}X{]} 的可逆元素群 U 因此恰好是 K*。于是两个相伴多项式仅相差一个非零常数因子；特别地，每个相伴多项式类包含唯一一个首一多项式。乘法群 K(X)* 的由首一多项式生成的子群因此从每个相伴有理函数类中恰好包含一个元素，从而与群

\[
\mathcal {P} ^ {*} = \mathrm{K} (\mathrm{X}) ^ {*} / \mathrm{U}
\]

\(\mathrm{K}(\mathrm{X})\) 的主分式理想组成的群。特别地，每当提及域 \(\mathrm{K}(\mathrm{X})\) 中（关于环 \(\mathrm{K}[\mathrm{X}]\) ）的 gcd 或 lcm 时，通常理解为这些是首一多项式（或 0）的商；这一约定使我们能够谈论一族有理函数的 gcd 或 lcm。

K{[}X{]} 的不可约元素正是通常意义下的不可约多项式（IV，第 13 页，定义 2），而首一不可约多项式的集合构成了 K{[}X{]} 的不可约元素的一个代表系.

一次多项式总是不可约的。若 K 是代数闭域，则其逆命题成立（V，第 19 页，命题 1）；因此在这种情况下，每个多项式 \(p(\mathbf{X})\) （ \(n\) 次的，属于 \(\mathbf{K}[\mathbf{X}]\) ）都可以唯一地（不计因子的顺序）写成如下形式

\[
p (\mathbf {X}) = c \left(\mathbf {X} - a _ {1}\right) \left(\mathbf {X} - a _ {2}\right) \dots \left(\mathbf {X} - a _ {n}\right)
\]

，其中 \(c\) 与诸 \(a_{i}\) 都是 \(\mathbf{K}\) 的元素.

命题6.——对于任意域 K，K{[}X{]} 中首一不可约多项式的集合是无限的。

事实上，给定任意非空有限族 \((p_{i})\) ( \(1 \leqslant i \leqslant n\) ) 由不同的首一不可约多项式组成，多项式 \(\left(\prod_{i=1}^{n} p_{i}\right) + 1\) 不可逆，且该多项式的任何首一不可约因子 q 必然与所有 \(p_{i}\) 不同，否则它会整除 1。

\section{主理想整环上的挠模}

\subsection{环的乘积上的模}

设 \(\mathbf{A}\) 是一个环，且 \((\mathfrak{b}_i)_{i\in I}\) 是 \(\mathbf{A}\) 的一个直和分解，即（I，第 110 页，定义 7）\(\mathbf{A}\) 的双边理想组成的一个有限族，使得从 \(\mathbf{A}\) 到诸 \(\mathbf{A} / \mathfrak{b}_i\) 的乘积的自然同态是双射的。根据前引文献命题 10，存在一个由中心幂等元组成的族 \((e_i)_{i\in I}\) （ \(\mathbf{A}\) 的），使得 \(\mathfrak{b}_i = \mathbf{A}(1 - e_i)\) ， \(\sum_{i\in I}e_i = 1\) ，且 \(e_i e_j = 0\) 对 \(i\neq j\) 成立.

对于每个左 A-模 M，令 \(M_{i}\) 表示由 \(m \in M\) （满足 \(b_{i}m = 0\) ）组成的集合；由于 \(b_{i}\) 是一个双边理想，这是 M 的一个子模；此外，如果 \(a, b \in A\) 且 \(a - b \in b_{i}\) ，则位似变换 \(a_{M_{i}}\) 与 \(b_{M_{i}}\) 重合；因此存在唯一的 \((\mathbf{A}/\mathbf{b}_{i})\) -模结构（在 \(M_{i}\) 上），使得 \(M_{i}\) 的 A-模结构经由同态 \(A \to A/b_{i}\) 诱导而来.

命题1. --- A-模 M 是其子模 \(M_{i}\) 的直和.

设 \(p_i: \mathbf{M} \to \mathbf{M}\) 表示位似变换 \(m \mapsto e_i m\) ；映射 \(p_i\) 是 A-线性的，因为 \(e_i\) 是中心的；由于 \(e_i^2 = e_i\) ， \(\sum_{i \in I} e_i = 1\) ，且 \(e_i e_j = 0\) 对 \(i \neq j\) 成立，我们有

\[
p _ {i} \circ p _ {i} = p _ {i}, \quad \sum_ {i \in I} p _ {i} = 1 _ {\mathrm{M}}, \quad p _ {i} \circ p _ {j} = 0 \quad \text { for } \quad i \neq j,
\]

，而诸 \(p_i\) 构成一个正交投影算子族，其和为恒等算子（II，第 209 页，定义 7）。根据前引文献命题 12，模 \(\mathbf{M}\) 是子模 \(p_i(\mathbf{M}) = e_i\mathbf{M}\) 的直和。此外 \(e_i\mathbf{M}\) 被 \(\mathfrak{b}_i = \mathbf{A}(1 - e_i)\) 零化；如果 \(i \neq j\) 且 \(m \in \mathbf{M}\) ，则 \((1 - e_i)e_jm = e_jm\) ，因此 \(e_j\mathbf{M}\) 中没有任何非零元素被 \(1 - e_i\) 零化，更不用说被 \(\mathfrak{b}_i\) 零化了。由此可知， \(e_i\mathbf{M} = \mathbf{M}_i\) ，命题得证。

注记. --- 1) 反之，设 \(\mathbf{M}_i'\) 是一个 \((\mathbf{A} / \mathfrak{b}_i)\) -模（对每个 \(i\) ），并考虑 A-模 \(\mathbf{M}\) ，即诸 A-模 \(\mathbf{M}_i'\) 的直和；则上述构造的子模 \(\mathbf{M}_i\) 与 \(\mathbf{M}_i'\) 重合（只需注意若 \(i \neq j\) ，则 \(\mathbf{M}_j'\) 中没有非零元素被 \(\mathfrak{b}_i\) 零化，因为 \(\mathfrak{b}_i + \mathfrak{b}_j = \mathbf{A}\) ). 因此，粗略地说，考虑一个 A-模 \(\mathbf{M}\) ，或模的一个族 \((\mathbf{M}_i)\) （在环 \(\mathbf{A} / \mathfrak{b}_i = \mathbf{A}_i\) 上的），是一回事.

\begin{enumerate}
\def\labelenumi{\arabic{enumi})}
\setcounter{enumi}{1}
\item
  由上述证明，M 到各分量的投影算子 \(M_{i}\) 是位似变换。
\item
  A-模 \(\mathbf{M}\) 是循环的当且仅当每个 \(\mathbf{M}_i\) 是循环的：如果 \(\mathbf{M} = \mathbf{A}m\) ，则 \(\mathbf{M}_i = \mathbf{A}_i e_i m\) ；反之，如果 \(\mathbf{M}_i = \mathbf{A}_i m_i\) 且 \(m = \sum_{i \in I} m_i\) ，则 \(\mathbf{M} = \mathbf{A}m\) ；事实上，如果
\end{enumerate}

\(n \in M\) 投影到 \(a_i m_i\) （对每个 \(i\) ），且若 \(a \in A\) 与 \(a_i\) 模 \(b_i\) 同余（对每个 \(i\) ），则 \(am\) 与 \(n\) 在每个 \(M_i\) 中具有相同的像，因此重合。

\begin{enumerate}
\def\labelenumi{\arabic{enumi})}
\setcounter{enumi}{3}
\tightlist
\item
  设 \(\mathbf{M}\) 与 \(\mathbf{N}\) 是两个 A-模，其分量分别为 \((\mathbf{M}_i)\) 与 \((\mathbf{N}_i)\) 。设 \(u \in \operatorname{Hom}_{\mathbf{A}}(\mathbf{M}, \mathbf{N})\) 是从 \(\mathbf{M}\) 到 \(\mathbf{N}\) 的 A-线性映射；则对所有 \(i\) 且对所有 \(m \in \mathbf{M}_i\) ，我们有 \(u(m) \in \mathbf{N}_i\) ，所以 \(u\) 诱导一个 \(\mathbf{A}_i\) -线性映射 \(u_i \in \operatorname{Hom}_{\mathbf{A}_i}(\mathbf{M}_i, \mathbf{N}_i)\) . 容易验证映射 \(u \mapsto (u_i)\) 是 \(\mathbf{Z}\) -模（当 \(\mathbf{A}\) 交换时则为 A-模）的同构
\end{enumerate}

\[
\operatorname{Hom} _ {\mathbf {A}} (\mathbf {M}, \mathbf {N}) \rightarrow \prod_ {i \in I} \operatorname{Hom} _ {\mathbf {A} _ {i}} (\mathbf {M} _ {i}, \mathbf {N} _ {i}).
\]

\subsection{主理想整环上挠模的典范分解}

设 M 是交换环 A 上的一个模。对于每个 \(\alpha\in A\) ，令 \(\mathbf{M}(\alpha)\) 表示 M 的自同态 \(x\mapsto\alpha x\) 的核。如果 \(\alpha\) 与 \(\beta\) 是 A 的两个元素，使得 \(\alpha\) 整除 \(\beta\) ，那么显然 \(\mathbf{M}(\alpha)\subset\mathbf{M}(\beta)\) 。特别地，当 n 遍历 \(\geqslant1\) 的有理整数集时，子模 \(\mathbf{M}(\alpha^{n})\) 构成一个递增序列；并集 \(M_{\alpha}\) （诸 \(\mathbf{M}(\alpha^{n})\) 的）因此是 M 的一个子模，由 M 中被 \(\alpha\) 的某个幂零化的元素组成。对于 M 的每个子模 N，显然有 \(N_{\alpha}=N\cap M_{\alpha}\) .

定义 1. --- 设 \(\pi\) 是主理想整环 A 的一个不可约元素；一个 A-模 M 被称为 \(\pi\) -准素的，如果对所有 \(x \in \mathbf{M}\) ，存在一个整数 \(n \geqslant 1\) 使得 \(\pi^n x = 0\) （换言之，如果 M 等于子模 \(\mathbf{M}_{\pi}\) ).

显然，每个形如 \(\mathbf{A}/(\pi^{s})\) 的循环模都是 \(\pi\) -准素的。对于任意一个 A-模 M，子模 \(M_{\pi}\) 是 \(\pi\) -准素的。

引理 1. --- 设 M 是主理想整环 A 上的一个模；对所有 \(\alpha \in \mathbf{A}\) （满足 \(\alpha \neq 0\) ），设 \(\alpha = \varepsilon \sum_{i=1}^{r} \pi_i^{n(i)}\) 是 \(\alpha\) 分解为不可约因子的一个分解（VII，第 4 页）。子模 \(\mathbf{N} = \mathbf{M}(\alpha)\) （由 M 中被 \(\alpha\) 零化的元素组成）是诸子模 \(\mathbf{M}(\pi_i^{n(i)})\) 的直和，且把每个 \(x \in \mathbf{M}(\alpha)\) 映到它在 \(\mathbf{M}(\pi_i^{n(i)})\) 中的分量的映射具有形式 \(x \mapsto \gamma_i x\) ( \(\gamma_i \in \mathbf{A}\) ）。此外

\[
\mathbf {M} \left(\pi_ {i} ^ {n (i)}\right) = \mathbf {N} \cap \mathbf {M} _ {\pi_ {i}} = \mathbf {N} _ {\pi_ {i}}.
\]

首先注意，N 被 \(\alpha\) 零化，因此具有自然的 \(\mathbf{A}/(\alpha)\) -模结构。根据 VII 第 3 页命题 4，从 \(\mathbf{A}/(\alpha)\) 到环的乘积 \(\mathbf{A}/(\pi_{i}^{n(i)})\) 的典范同态是环同构；现在应用 VII 第 6 页命题 1，我们推出 N 是 \(\mathbf{M}(\pi_{i}^{n(i)})\) 的直和；该分解的投影算子是位似变换，依据 VII 第 7 页注记 2。包含关系 \(\mathbf{M}(\pi_{i}^{n(i)})\subset\mathbf{M}(\alpha)\cap\mathbf{M}_{\pi_{i}}\) 是显然的；反之，设 \(x\in\mathbf{M}(\alpha)\cap\mathbf{M}_{\pi_{i}}\) 于是存在一个幂 \(\pi_{i}^{s}\) （ \(\pi_{i}\) 的幂）零化 x；我们可以假设 \(s\geqslant n(i)\) ；根据贝祖恒等式，存在 \(\lambda,\mu\in\mathbf{A}\) 使得 \(\pi_{i}^{n(i)}=\lambda\pi_{i}^{s}+\mu\alpha\) ，所以 \(\pi_{i}^{n(i)}x=0\) ，最终得 \(x\in\mathbf{M}(\pi_{i}^{n(i)})\) .

引理 2. —— 设 M 是整环 A 上的挠模（II，第 313 页）。对于 M 中元素的每个有限族 \((x_{i})_{1 \leq i \leq n}\) ，存在 A 中的一个元素 \(\gamma \neq 0\) 使得 \(x_{i}\) 都属于 \(\mathrm{M}(\gamma)\) .

事实上，对于每个指标 \(i\) ，在 A 中存在一个元素 \(\alpha_{i} \neq 0\) 零化 \(x_{i}\) ，且元素 \(\gamma = \prod_{i=1}^{n} \alpha_{i}\) 即符合要求。

定理1. --- 设 M 是主理想整环 A 上的挠模；对 A 的每个不可约元素 \(\pi\) ，令 \(M_{\pi}\) 为 M 的由被 \(\pi\) 的某个幂零化的元素组成的子模。若 P 是 A 的不可约元素的一个代表系，则 M 是其子模 \(M_{\pi}\) （ \(\pi \in P\) ）的直和.

每个元素 \(x \in M\) 都属于某个子模 \(\mathbf{M}(\alpha)\) （ \(\alpha \neq 0\) ），因此根据引理 1，它是有限个元素的和，其中每个元素都属于某个子模 \(M_{\pi}\) 。另一方面，如果 \(\sum_{\pi \in P} x_{\pi} = \sum_{\pi \in P} y_{\pi}\) ，其中 \(x_{\pi}, y_{\pi} \in M_{\pi}\) 对所有 \(\pi \in P\) 成立，且其中 \(x_{\pi}\) 与 \(y_{\pi}\) 除有限多个外均为零，则引理 2 表明存在 A 中的 \(\gamma \neq 0\) ，使得所有 \(x_{\pi}\) 与 \(y_{\pi}\) 都属于同一个子模 \(\mathbf{M}(\gamma)\) ；将引理 1 应用于 \(\mathbf{M}(\gamma)\) 表明 \(x_{\pi} = y_{\pi}\) 对所有 \(\pi \in P\) 成立，这就完成了证明。

显然，如果 \(\pi\) 与 \(\pi'\) 是两个相伴的不可约元素，则 \(M_{\pi} = M_{\pi'}\) ；因此，对于给定的模 M，子模 \(M_{\pi}\) 仅取决于 A 的理想 ( \(\pi\) )；它称为模 M 的 \(\pi\) -准素分量，而将 M 分解为诸 \(M_{\pi}\) 的直和称为 M 作为其 \(\pi\) -准素分量直和的典范分解。

推论1. --- 挠模 M 的每个子模 N 都是其子模 \(N \cap M_{\pi}\) 的直和.

这由以下事实得出： \(\mathbf{N} \cap \mathbf{M}_{\pi}\) 是 \(\pi\) -准素分量 \(\mathbf{N}_{\pi}\) （ \(\mathbf{N}\) 的）。

推论2.——挠 A-模 M 的子模 N 是直和因子，当且仅当 \(N_{\pi}\) 是 \(M_{\pi}\) （对 A 的每个不可约元素 \(\pi\) ）的直和因子。

确实，如果 \(\mathbf{N}\) 与 \(\mathbf{N}'\) 是 \(\mathbf{M}\) 的两个子模，则 \(\mathbf{M} = \mathbf{N} \oplus \mathbf{N}'\) 当且仅当 \(\mathbf{M}_{\pi} = \mathbf{N}_{\pi} \oplus \mathbf{N}_{\pi}'\) 对每个不可约元素 \(\pi\) （ \(\mathbf{A}\) 的）成立（推论 1）。

推论3.——设 N 是挠 A-模 M 的一个子模。如果对于 A 的每个不可约元素 \(\pi\) ，要么 \(N_{\pi} = 0\) ，要么 \((M/N)_{\pi} = 0\) ，则 N 是 M 的一个直因子。

事实上，条件 \((\mathbf{M}/\mathbf{N})_{\pi}=0\) 蕴含 \(N_{\pi}=M_{\pi}\) ，于是推论 2 适用。

一个 A-模 M 称为半单的，如果 M 的每个子模都是直因子（参见 A，VIII，§3）。

推论4. --- 设 A 是一个不是域的主理想整环，并设 M 是一个 A-模。则 M 是半单的，当且仅当 M 是挠模，且 \(\mathbf{M}_{\pi} = \mathbf{M}(\pi)\) 对 A 的每个不可约元素 \(\pi\) 成立。

首先假设 M 是半单的；令 \(x \in M\) 并且令 \(\pi\) 是 A 的一个不可约元素。若 N 是 \(A\pi x\) 在 M 中的一个补，那么我们可以写出 \(x = \alpha\pi x + y\) ，其中 \(\alpha \in A\) 且 \(y \in N\) ；但这意味着 \(y = (1 - \alpha\pi)x\) ，所以

\[
\pi (1 - \alpha \pi) x \in A \pi x \cap N = 0.
\]

首先由此得出 M 是一个挠模；如果此外 \(x \in M_{\pi}\) ，则 \(\pi(1 - \alpha\pi)x = 0\) ，从而 \(\pi x = \alpha\pi^{2}x = \alpha^{2}\pi^{3}x = \cdots = \alpha^{n}\pi^{n+1}x\) 为零，且 \(\mathbf{M}_{\pi} = \mathbf{M}(\pi)\) .

反之，由推论 2 可知，只需证明被一个不可约元素 \(\pi\) 零化的 A-模 M 是半单的；但这是显然的，因为此时 M 具有域 \(\mathrm{A}/(\pi)\) 上向量空间的自然结构，并且在此结构下，M 的子模恰好就是向量子空间。

注记 1. --- 显然，每个 \(\neq 0\) 元素（ \(\pi\) -准素模中的）的零化子都具有形式 \(A\pi^k (k > 0\) 为一个整数），因为它是包含 \(\pi\) 的一个幂的主理想。设 \(x\) 为 \(M\) 的一个元素；对每个 \(\pi \in P\) ，设 \(x_{\pi}\) 为 \(x\) 在 \(M_{\pi}\) 中的分量； \(x\) 的零化子是诸非零 \(x_{\pi}\) 的零化子的最小公倍元，但根据上述内容，在这种情形下它等于诸非零 \(x_{\pi}\) 的零化子的乘积（VI，第 16 页，命题 12（DIV））。

命题2. --- 若 M 是主理想整环 A 上的有限生成挠模，则 M 的 \(\pi\) -准素分量除有限个外均为零，且 M 在这些分量 \(M_{\pi}\) 上的投影算子是位似变换。

这由引理1直接得出，因为根据引理2，存在 \(\alpha \neq 0\) 在 A 中使得 \(\mathbf{M} = \mathbf{M}(\alpha)\) .

注2. --- 由 VII 第 7 页注记 3，一个有限生成的挠 A-模是循环的，当且仅当其每个 \(\pi\) -准素分量都是循环的。

定理1和命题2适用的一个重要特殊情形是当A=Z时；一个Z-模就是一个阿贝尔群。一个阿贝尔群被称为挠群，如果它是一个挠Z-模，也就是说它的所有元素都有有限阶。此时，取 P 为 \textgreater{} 0 的素数集合；对于每个素数 p \textgreater{} 0，一个（阿贝尔）群被称为 p-挠群，如果其所有元素的阶都是 p 的幂。在这个术语下，定理 1 表明每个挠阿贝尔群都是 p-挠群的直和。在有限群的情形下，这也由I，第80页，定理4得出。

\subsection{应用：I. 有理数及一元有理函数的典范分解}

定理2.——设A为主理想整环，K为其分式域，P为A中不可约元素的代表元系。给定一个元素 \(x \in K\) ，存在 P 的一个有限子集 H、元素 \(a_{0} \in A\) 与 \(a_{p} \in A\) （在 A 中不被 p 整除， \(p \in H\) ）、以及整数 \(s(p) > 0\) ( \(p \in H\) ）使得

\[
x = a _ {0} + \sum_ {p \in \mathrm{H}} a _ {p} p ^ {- s (p)},\tag{1}
\]

其中H和 \(s(p)\) 由这些条件唯一确定。

此外，如果 \(R_{p}\) 表示 A 的一个恰好包含 A 中每个模 p 剩余类中一个元素的子集（ \(p \in P\) ），则每个 \(x \in K\) 可以唯一地表示为

\[
x = a + \sum_ {p \in P} \left(\sum_ {h = 1} ^ {\infty} r _ {p h} p ^ {- h}\right)\tag{2}
\]

，其中 \(a \in \mathbf{A}\) ，且 \(r_{ph} \in \mathbf{R}_p\) 对所有 \(h\) 与 \(p\) 成立，并且除有限多个外 \(r_{ph}\) 均为零。

将 K 视为 A-模；则 A 是由 1 生成的子模。商模 K/A 是 K 关于等价关系 \(x' - x \in \mathbf{A}\) 的商，在 VI 第 6 页的记号中也写作 \(x \equiv x' \pmod{1}\) ；设 \(f\) 表示从 K 到 M = K/A 上的自然同态。

商模 M 是一个挠模，因为 M 的每个元素都具有形式 \(f(a/b)\) ( \(a \in A\) , \(b \in A\) , \(b \neq 0\) ），从而 \(bf(a/b) = f(a) = 0\) 。因此，VII 第 8 页的定理 1 适用。设 \(M_{p}\) ( \(p \in P\) ) 为 M 中被 p 的幂零化的元素构成的子模；则 \(f^{-1}(M_{p})\) 是子环 \(A_{p}\) （由 K 中形如 \(ap^{-n}\) 的元素组成，其中 \(a \in A\) ， \(n \geqslant 0\) 为整数）。模 M 是这些 \(M_{p}\) 的直和，因此每个 \(x \in K\) 在模 1 下同余于这些 \(A_{p}\) 之和中的一个元素；换言之，公式 (1) 成立，其中 \(s(p)\) 为 \textgreater0 的整数，且诸 \(a_{p}\) 是 A 中有限多个元素，使得 \(a_{p}\) 不是 p 的倍数。

我们现在证明这些关于 \(s(p)\) 与 \(a_{p}\) 的条件完全确定了 H 和 \(s(p)\) 。事实上，H 就是 \(p \in \mathbf{P}\) 中使得 \(f(x)\) 在 \(\mathbf{M}_{p}\) 中的分量 \(\neq 0\) 的那些 p 的集合。另一方面，如果 \(s\) 与 \(s'\) 是两个 \(>0\) 的整数，满足 \(s \geqslant s'\) ，并且 \(a\) 与 \(a'\) 是 \(A\) 的不被 \(p\) 整除的元素，满足 \(ap^{-s} \equiv a'p^{-s'} \pmod{1}\) ，那么我们推断出 \(a \equiv a'p^{s - s'} \pmod{p^s}\) ；如果 \(s > s'\) ，则 \(a \equiv 0 \pmod{p}\) ，这与假设相矛盾。这个论证还表明每个 \(a_p\) 在模 \(p^{s(p)}\) 下是唯一确定的.

为了完成证明，我们首先注意到，在每个模 \(p^s\) 的剩余类（在 \(\mathbf{A}\) 中）里，存在唯一一个形如 \(\sum_{h=0}^{s-1} r_h p^h\) 的元素，其中 \(r_h \in \mathbb{R}_p\) ， \(0 \leqslant h \leqslant s-1\) 。事实上，我们对 \(s\) 用归纳法（该性质由 \(\mathbb{R}_p\) 的定义在 \(s=1\) 时得出）：设 \(x \in \mathbf{A}\) ；根据假设，存在唯一一个形如 \(\sum_{h=0}^{s-2} r_h p^h\) ( \(r_h \in \mathbb{R}_p\) ) 的元素，属于 \(x \mod p^{s-1}\) 的剩余类；则差 \(x - \sum_{h=0}^{s-2} r_h p^h\) 形如 \(ap^{s-1}\) ，即 \(p^{s-1}\) 的一个倍数；现在存在唯一的元素 \(r_{s-1}\) （ \(\mathbb{R}_p\) 的）使得 \(a \equiv r_{s-1} (\bmod p)\) ；由此 \(x \equiv \sum_{h=0}^{s-1} r_h p^h (\bmod p^s)\) 。为了得到公式 (2)，现在只需将此事实应用于公式 (1) 中的每个 \(a_p\) 。鉴于上述，唯一性是显然的。

以下是定理2最重要的应用：

I. 环 A 是有理整数环 Z，且 K = Q。设 P 为 \textgreater{} 0 的素数集合，且对于每个 \(p \in P\) ，令 \(R_{p}\) 为 Z 中的区间 \((0, p - 1)\) 。因此我们有典范分解

\[
x = a + \sum_ {p \in P} \left(\sum_ {h = 1} ^ {\infty} e _ {p h} p ^ {- h}\right)
\]

，其中 \(a \in \mathbf{Z}\) , \(e_{ph} \in \mathbf{Z}\) ，且 \(0 \leqslant e_{ph} \leqslant p - 1\) .

\begin{enumerate}
\def\labelenumi{\Roman{enumi}.}
\setcounter{enumi}{1}
\tightlist
\item
  环 A 是交换域 E 上一个未定元的多项式环 E{[}X{]}，且 \(\mathbf{K} = \mathbf{E}(\mathbf{X})\) 。设 P 为 E{[}X{]} 中首一不可约多项式的集合（VII，第 5 页）。对于 \(p \in P\) ，根据多项式的欧几里得除法（IV，第 10 页），我们可以取 \(R_{p}\) 为次数严格小于 p 的次数的多项式集合。因此我们得到有理函数 \(r(\mathbf{X}) \in \mathbf{E}(\mathbf{X})\) 的（称为典范分解的）分解：
\end{enumerate}

\[
r (\mathrm{X}) = a (\mathrm{X}) + \sum_ {p \in \mathrm{P}} \left(\sum_ {h = 1} ^ {\infty} v _ {p h} (\mathrm{X}) \cdot p (\mathrm{X}) ^ {- h}\right)
\]

，其中 \(a(\mathbf{X})\) 是一个多项式，且 \(v_{ph}(\mathbf{X})\) 是次数严格小于 \(p(\mathbf{X})\) 的次数的多项式（对所有 \(p\) 与 \(h\) ）。特别地，如果域 \(\mathbf{E}\) 是代数闭的，则诸 \(p(\mathbf{X})\) 具有形式 \(\mathbf{X} - \alpha\) （ \(\alpha \in \mathbf{E}\) ），从而诸 \(v_{ph}(\mathbf{X})\) 是常数。

因此我们可以说，域 E 上的向量空间 \(E(X)\) 具有一个基，由单项式 \(X^{n}\) ( \(n \geqslant 0\) 为整数）以及形如 \(\mathbf{X}^{m}/(p(\mathbf{X}))^{h}\) 的有理函数组成，其中 \(p \in P\) ，h 和 m 是满足 \(h \geqslant 1\) 与 \(0 \leqslant m < \deg(p)\) 的整数.

\subsection{应用：II. 整数模某数的单位乘法群}

设 \(a\) 为一个有理整数 \(> 1\) ，并设 \((\mathbf{Z} / a\mathbf{Z})^*\) 为环 \(\mathbf{Z} / a\mathbf{Z}\) 的可逆元素的乘法群。若 \(a = \prod_{i} p_i^{n(i)}\) 是 \(a\) 分解为

素因子的分解，则环 \(\mathbf{Z}/a\mathbf{Z}\) 同构于环 \(\mathbf{Z}/p_{i}^{n(i)}\mathbf{Z}\) 的乘积（VII，第 3 页，命题 4），且群 \((\mathbf{Z}/a\mathbf{Z})^{*}\) 同构于群 \((\mathbf{Z}/p_{i}^{n(i)}\mathbf{Z})^{*}\) 的乘积。因此我们归结为研究群 \((\mathbf{Z}/p^{n}\mathbf{Z})^{*}\) ，其中 \(p\) 是一个素数；回顾（V，第 80 页）， \(\varphi(p^{n})\) （ \((\mathbf{Z}/p^{n}\mathbf{Z})^{*}\) 的阶）是 \(p^{n}-p^{n-1}=p^{n-1}(p-1)\) .

首先假设 \(p > 2\) ；自然同态 \(\mathbf{Z} / p^n\mathbf{Z} \to \mathbf{Z} / p\mathbf{Z}\) 限制为从 \((\mathbf{Z} / p^n\mathbf{Z})^*\) 到 \((\mathbf{Z} / p\mathbf{Z})^*\) 上的群同态，其核记为 \(\mathrm{U}(p^n)\) ；由 VII 第 3 页命题 3，模 \(p^n\) 的剩余类（整数 \(m\) 的）可逆当且仅当 \(m\) 与 \(p\) 互素，即当且仅当 \(m \mod p\) 的剩余类可逆。由此可知 \(\mathrm{U}(p^n)\) 由所有模 \(p^n\) 与 \(1 \mod p\) 同余的整数的剩余类组成，因此有 \(p^{n-1}\) 个元素，并且存在一个正合序列

\[
\{1 \} \rightarrow \mathrm{U} (p ^ {n}) \rightarrow (\mathbf {Z} / p ^ {n} \mathbf {Z}) ^ {*} \rightarrow (\mathbf {Z} / p \mathbf {Z}) ^ {*} \rightarrow \{1 \}.\tag{3}
\]

类似地，对于 \(n \geqslant 2\) ，令 \(U(2^n)\) 表示从 \((\mathbf{Z}/2^n\mathbf{Z})^*\) 到 \((\mathbf{Z}/4\mathbf{Z})^*\) 的自然同态的核；这是一个 \(2^{n-2}\) 阶的群，由所有模 \(2^n\) 与 1 模 4 同余的整数的剩余类组成，并且存在一个正合序列

\[
\{1 \} \rightarrow \mathrm{U} (2 ^ {n}) \rightarrow (\mathbf {Z} / 2 ^ {n} \mathbf {Z}) ^ {*} \rightarrow (\mathbf {Z} / 4 \mathbf {Z}) ^ {*} \rightarrow \{1 \}.\tag{4}
\]

引理 3. --- 设 \(x, y, k\) 为整数，满足 \(k \geqslant 0\) ，并设 \(p > 2\) 为素数。若 \(x \equiv 1 + py \mod p^2\) ，则 \(x^{p^k} \equiv 1 + p^{k+1}y \mod p^{k+2}\) 。若 \(x \equiv 1 + 4y \mod 8\) ，则 \(x^{2^k} \equiv 1 + 2^{k+2}y \mod 2^{k+3}\) .

为证明第一个断言，只需证明：若 \(k \geqslant 1\) 且 \(x \equiv 1 + p^{k}y \mod p^{k+1}\) ，则 \(x^{p} \equiv 1 + p^{k+1}y \mod p^{k+2}\) ，然后对整数 k 作归纳论证。对所有 \(a \in Z\) 与 \(k \geqslant 1\) ，显然有

\[
(1 + p ^ {k} a) ^ {p} \equiv 1 + p ^ {k + 1} a \mod p ^ {k + 2},
\]

因此

\[
\begin{array}{r l} (1 + p ^ {k} y + p ^ {k + 1} z) ^ {p} & = (1 + p ^ {k} (y + p z)) ^ {p} \equiv \\ & \equiv 1 + p ^ {k + 1} (y + p z) \equiv 1 + p ^ {k + 1} y \mod p ^ {k + 2} \end{array}
\]

类似地，对于 \(k \geqslant 1\) 我们有

\[
(1 + 2 ^ {k + 1} a) ^ {2} \equiv 1 + 2 ^ {k + 2} a \mod 2 ^ {k + 3},
\]

所以

\[
(1 + 2 ^ {k + 1} y + 2 ^ {k + 2} z) ^ {2} \equiv 1 + 2 ^ {k + 2} y \mod 2 ^ {k + 3},
\]

由此根据对k的归纳即得第二个断言。

命题 3. --- 设 p \textgreater{} 2 为素数， n \textgreater{} 0 为整数；则群 U(p\^{}n) 是 p\^{}\{n-1\} 阶的循环群；若 n ≥ 2，则与 1 模 p 同余的整数 x 的模 p\^{}n 剩余类是 U(p\^{}n) 的生成元，当且仅当 x 不同余于 1 模 p\^{}2。设 m \textgreater{} 1 为整数；则群 U(2\^{}m) 是 2\^{}\{m-2\} 阶的循环群；若 m ≥ 3，则与 1 模 4 同余的整数 x 的模 2\^{}m 剩余类是 U(2\^{}m) 的生成元，当且仅当 x 不同余于 1 模 8。

由于 \(\mathrm{U}(p^{n})\) 的阶为 \(p^{n-1}\) ， \(\mathrm{U}(p^{n})\) 的每个元素 u 的阶都是 p 的幂，且 u 是 \(\mathrm{U}(p^{n})\) 的生成元当且仅当 \(u^{p^{n-2}} \neq 1\) 。现在若 u 是 \(x = 1 + py\) 的类，则 \(u^{p^{n-2}}\) 是 \(1 + p^{n-1}y\) 的类（由引理 3），因此 u 生成 \(\mathrm{U}(p^{n})\) 当且仅当 \(y \neq 0 \mod p\) ，换言之 \(x \neq 1 \mod p^{2}\) 。例如，类 \(1 + p\) 生成 \(\mathrm{U}(p^{n})\) 。类似地，x 模 \(2^{n}\) 的类 u 生成 \(\mathrm{U}(2^{n})\) 当且仅当 \(u^{2^{n-3}} \neq 1\) ，这意味着 x 不同余于 1 模 8（由引理 3）；x = 5 便满足这一点。

引理4. --- 设 A 是一个主理想整环，且设 \(0 \rightarrow N \rightarrow M \rightarrow P \rightarrow 0\) 是 A-模的一个正合序列。假设存在互素元素 a, \(b \in A\) 使得aN = 0且bP = 0。则正合序列分裂。若此外N和P均为循环群，则M为循环群。

模 M 是挠模，因为 abM = 0。第一个断言由 VII 第 9 页推论 3 得出。若 N 和 P 是循环模，则它们是有限生成的，因此 M 也是有限生成的（II 第 17 页推论 5）；由于 M 的每个 p-准素分量同构于 N 或 P 的某个 p-准素分量，由 VII 第 10 页注记 2 可知 M 是循环模。

定理3. --- 若 \(a = \prod_{i} p_{i}^{n(i)}\) 是整数的质因数分解

为 \textgreater{} 1，则环 Z/aZ 的可逆元素群 \((\mathbf{Z}/a\mathbf{Z})^{*}\) 同构于诸群 \((\mathbf{Z}/p_{i}^{n(i)}\mathbf{Z})^{*}\) 的乘积。如果 p \textgreater{} 2 是一个素数且 \(n \geqslant 1\) 是一个整数，则群 \((\mathbf{Z}/p^{n}\mathbf{Z})^{*}\) 是 \(p^{n-1}(p - 1)\) 阶的循环群。群 \((\mathbf{Z}/2\mathbf{Z})^{*}\) 是平凡的；对于 \(n \geqslant 2\) ，群 \((\mathbf{Z}/2^{n}\mathbf{Z})^{*}\) 是 \(2^{n-2}\) 阶循环群与 2 阶循环群的直积，前者由 5 模 \(2^{n}\) 的剩余类生成，后者由 1 和 -1 模 \(2^{n}\) 的剩余类组成。

\(p^{n-1}\) （ \(U(p^n)\) 的阶）与 \(p-1\) （ \((\mathbf{Z}/p\mathbf{Z})^*\) 的阶）互素；由于 \(U(p^n)\) 与 \((\mathbf{Z}/p\mathbf{Z})^*\) 都是循环群（命题 3 及 V 第 78 页引理 1），群 \((\mathbf{Z}/p^n\mathbf{Z})^*\) 是循环群（将引理 4 应用于正合序列 (3)）。若 \(n \geqslant 2\) ，则同态 \(v: (\mathbf{Z}/2^n\mathbf{Z})^* \to (\mathbf{Z}/4\mathbf{Z})^*\) 在子群 \(\{1, -1\}\) 上的限制是双射；因此群 \((\mathbf{Z}/2^n\mathbf{Z})^*\) 是该子群与核 \(U(2^n)\) （ \(v\) 的）的直积；结论由命题 3 得出。

注记. --- 设 p \textgreater{} 2 为素数，且设 x 为与 1 模 p 同余而不同余于 1 模 \(p^{2}\) 的整数；存在一个正合序列

\[
\{0 \} \rightarrow \mathbf {Z} / p ^ {n - 1} \mathbf {Z} \xrightarrow {u} (\mathbf {Z} / p ^ {n} \mathbf {Z}) ^ {*} \xrightarrow {v} (\mathbf {Z} / p \mathbf {Z}) ^ {*} \rightarrow \{1 \}\tag{5}
\]

的群，其中 v 是自然投影，而 u 在商上由映射 \(r \mapsto x^{r}\) 诱导。设 \(Z_{p}\) 为 p 进整数环（V，第 96 页），并令 x 为 \(Z_{p}\) 中满足 \(x - 1 \in pZ_{p}\) 与 \(x - 1 \notin p^{2}Z_{p}\) 的元素；通过取逆向极限，正合序列 (5) 诱导出一个正合序列

\[
\{0 \} \rightarrow \mathbf {Z} _ {p} \stackrel {u} {\rightarrow} \mathbf {Z} _ {p} ^ {*} \stackrel {v} {\rightarrow} (\mathbf {Z} / p \mathbf {Z}) ^ {*} \rightarrow \{1 \}
\]

，其中 \(v\) 是自然映射，且连续映射 \(u\) 延拓了映射 \(n \mapsto x^n\) ( \(n \in \mathbf{Z}\) ）。我们常设 \(x^n = u(x)\) ，其中 \(n \in \mathbf{Z}_p\) .

类似地，如果 \(x \in \mathbf{Z}_2\) 满足 \(x - 1 \in 4\mathbf{Z}_2\) 且 \(x - 1 \notin 8\mathbf{Z}_2\) ，则存在一个分裂正合序列

\[
\{0 \} \rightarrow \mathbf {Z} _ {2} \stackrel {{u}} {{\rightarrow}} \mathbf {Z} _ {2} ^ {*} \stackrel {{v}} {{\rightarrow}} (\mathbf {Z} / 4 \mathbf {Z}) ^ {*} \rightarrow \{1 \},
\]

，其中 \(u\) 是映射 \(n \mapsto x^n\) 的一个连续延拓.

\section{主理想整环上的自由模}

定理1. --- 设A是一个环，使得A的每个左理想作为A-模都是投射的（II，第231页，定义1）。则自由左A-模L的每个子模M都是与A的理想同构的模的直和。

设 \((e_{\iota})_{\iota \in I}\) 是 L 的一个基，并令 \(p_{\iota}\) 为对应于该基的坐标函数。选取 I 的一个良序（集合论，III，第 153 页），并令 \(L_{\iota}\) 表示由诸 \(e_{\lambda}\) （ \(\lambda \leqslant \iota\) ）生成的子模；令 \(M_{\iota} = M \cap L_{\iota}\) 。坐标函数 \(p_{\iota}\) 把 \(M_{\iota}\) 映到 A 的一个理想 \(\mathfrak{a}_{\iota}\) 上；由于 \(\mathfrak{a}_{\iota}\) 是投射 A-模，存在（II，第 231 页，命题 4）一个子模 \(N_{\iota}\) （ \(M_{\iota}\) 的），使得映射 \(x \mapsto p_{\iota}(x)\) （从 \(N_{\iota}\) 到 \(\mathfrak{a}_{\iota}\) ）是双射的。设 \(M_{\iota}'\) 为 L 中由诸 \(N_{\lambda}\) （ \(\lambda \leqslant \iota\) ）生成的子模；我们将证明 \(M_{\iota}' = M_{\iota}\) 对所有 \(\iota\) 成立，由此将得出 M 由族 \((N_{\iota})_{\iota \in I}\) 生成。事实上，假设 \(M_{\lambda}' = M_{\lambda}\) 对所有 \(\lambda < \iota\) 成立；则 \(p_{\iota}(x) \in \mathfrak{a}_{\iota}\) 对所有 \(x \in M_{\iota}\) 成立；因此存在 \(y \in N_{\iota}\) 使得 \(x - y\) 是有限多个元素 \(e_{\lambda}\) （ \(\lambda < \iota\) ）的线性组合；换言之， \(x - y\) 是 \(M_{\lambda}\) （对某个 \(\lambda < \iota\) ）的元素；归纳假设表明 \(x - y \in M_{\lambda}' \subset M_{\iota}'\) ，即 \(x \in M_{\iota}'\) ，从而 \(M_{\iota}' = M_{\iota}\) 。还需证明诸 \(N_{\iota}\) 之和是直和；现在假设存在一个线性关系 \(\sum_{\iota} a_{\iota} = 0\) ，其中

\(a_{\iota} \in \mathbf{N}_{\iota}\) ，其中诸 \(a_{\iota}\) （除有限多个外均为零）并非全为零。设 \(\mu\) 为最大的指标 \(\iota\) ，使得 \(a_{\iota} \neq 0\) ；由于 \(p_{\mu}(a_{\lambda}) = 0\) （对 \(\lambda < \mu\) ），我们有 \(p_{\mu}(a_{\mu}) = p_{\mu}\left(\sum a_{\iota}\right) = 0\) ，所以 \(a_{\mu} = 0\) ，这与 \(\mu\) 的选取相矛盾.

推论1. --- 若A的每个左理想都是投射的，则投射左A-模的每个子模都是投射的。

确实每个投射A-模都是自由A-模的子模（II，第231页，命题4），因此定理1适用。

推论2. --- 主理想整环上自由模的每个子模都是自由的。

这直接由定理1推出，因为主理想整环的每个理想都是自由的。

推论3. --- 主理想整环上的每个投射模都是自由的。

注记. --- 定理 1 的证明表明， \(\mathbf{A}^{(1)}\) 的每个子模都同构于一个直和 \(\oplus \mathfrak{a}_{\iota}\) ，其中每个 \(\mathfrak{a}_{\iota}\) 是 \(\mathbf{A}\) 的理想.

命题1. --- 若L是主理想整环A上的秩为n的自由模，则L的每个子模M都是秩≤n的自由模。

确实，由定理 1 的推论 2，M 是自由模，且由前面的注记或由下面的引理，其秩为 \(\leqslant n\) ：

引理 1. --- 设 L 是交换环 A 上的一个由 n 个元素生成的模，且 M 是 L 的一个自由子模；则 M 的秩 ≤ n。

首先假设 L 是自由的。设 i 表示 M 到 L 中的典范单射。由 III 第 520 页推论，同态 \(\bigwedge_{n+1}^{n+1} i: \bigwedge_{n+1}^{n+1} M \to \bigwedge_{n+1}^{n+1} L\) 是单射；由 III 第 511 页命题 6， \(\bigwedge_{n+1}^{n+1} L = \{0\}\) ，所以 \(\bigwedge_{n+1}^{n+1} M = \{0\}\) ；由此可知 M 的秩为 \(\leqslant n\) （III，第 518 页，推论 1）。现在考虑一般情形；模 L 是一个秩为 n 的自由模 \(L'\) 的商。存在 L 的一个与 M 同构的子模 \(M'\) （II，第 218 页，命题 21）。由论证的第一部分， \(M'\) 的秩为 \(\leqslant n\) ，结论随之得出。

推论. --- 设 E 是主理想整环 A 上的模，由 n 个元素生成。则 E 的任意子模 F 可由至多 n 个元素生成。

事实上，存在一个同态 \(f\) ，从 \(\mathbf{A}^n\) 映到 \(\mathbf{E}\) 上（II，第 216 页，推论 3），且 \(f^{-1}(\mathbf{F})\) （它是秩为 \(m \leqslant n\) 的自由模）由 \(n\) 个元素生成；这些元素在 \(f\) 下的像生成 \(\mathbf{F}\) .

\section{主理想整环上的有限生成模}

\subsection{循环模的直和}

设 A 为一个交换环。回顾（II，第 220 页，命题 22），一个循环 A-模同构于商模 A/a，其中 a 是 A 的一个理想。我们将在后面（第 4 节）看到，主理想整环上的每个有限生成模都是有限多个循环模的直和。

命题 1. --- 设 E 为交换环 A 上的一个模；假设 E 是 n 个循环模 \(A/\mathfrak{a}_{k}\) ( \(1 \leqslant k \leqslant n\) ) 的直和，其中 \(a_{k}\) 是 A 的理想；

那么，对于每个整数 \(p > 0\) ，A-模 \(\bigwedge^p \mathbf{E}\) 同构于诸模 \(\mathbf{A} / \mathfrak{a}_{\mathrm{H}}\) 的直和，其中对于每个 \(p\) -元素子集 \(\mathbf{H} = \{k_1, \dots, k_p\}\) （ \([1, n]\) 的），理想 \(\mathfrak{a}_{\mathrm{H}}\) 是 \(\sum_{j=1}^{p} \mathfrak{a}_{k_j}\) .

设 \(x_{k}\) 为 \(A / a_{k}\) 的典范生成元，即 A 的单位元的像，因此 E 是诸 \(Ax_{i}\) ( \(1 \leqslant i \leqslant n\) ) 的直和。于是我们知道（III，第 515 页，命题 10），外代数 \(\bigwedge_{n} E\) 作为 A-模同构于张量积 \(\otimes_{i=1}^{n} (\bigwedge(Ax_{i}))\) 。现在 \(\bigwedge(Ax_{i})\) 正是直和 \(A \oplus Ax_{i}\) ，因为

\[
\bigwedge^ {p}
\]

\(Ax_{i}\) 中任意两个元素的外积为零，因此 \(\Lambda E\) 是诸模 \(\mathbf{M}_{\mathrm{H}} = (\mathbf{A}x_{k_{1}}) \otimes \ldots \otimes (\mathbf{A}x_{k_{p}})\) 的直和，其中 \(H = \{k_{1}, \ldots, k_{p}\}\) 取遍 {[}1, n{]} 的 p 元子集（ \(k_{1} < \ldots < k_{p}\) ）；现在已知 \(M_{H}\) 同构于 \(A/a_{H}\) （II，第 257 页，推论 4），这就完成了证明。

我们现在看到，在命题 1 的记号下，如果诸理想 \(\mathfrak{a}_k\) 构成一个递增序列，则它们完全由模 E 确定。更精确地说：

命题 2. --- 设 A 是一个交换环，且 E 是 n 个循环模 \(A/a_{k}\) 的直和，其中 \(a_{k}\) 满足 \(a_{1} \subset a_{2} \subset \ldots \subset a_{n}\) 。那么，对于 \(1 \leqslant p \leqslant n\) ，理想 \(a_{p}\) 是 \(\bigwedge^{p} E\) 的零化子；如果 \(a_{n} \neq A\) ，则 \(\bigwedge^{p} E \neq 0\) （对 \(1 \leqslant p \leqslant n\) ），且 \(\bigwedge^{m} E = 0\) （对 m \textgreater{} n）.

确实，在命题 1 的记号中，我们有 \(\mathfrak{a}_{\mathrm{H}} = \mathfrak{a}_{s(\mathrm{H})}\) ，其中 \(s(\mathrm{H})\) 表示子集 H 的最大元素。由于 \(s(\mathrm{H}) \geqslant p\) 对每个 \(p\) -元素子集 H 成立，且 \(s(\mathrm{H}) = p\) 对 \(\mathrm{H} = \{1, \dots, p\}\) 成立，由此可得， \(\mathfrak{a}_p\) 是诸 \(\mathfrak{a}_{\mathrm{H}}\) 的交，其中 H 取遍 \(p\) -元素子集（ \([1, n]\) 的）；因此理想 \(\mathfrak{a}_p\) 确实是 \(\bigwedge^p E\) 的零化子（由命题 1）。

推论. --- 在命题 2 的记号下，若 \(\mathfrak{a}_n \neq \mathbf{A}\) ，且 E 也同构于 \(m\) 个循环模 \(\mathbf{A} / \mathfrak{a}_j'\) 的直和，满足 \(\mathfrak{a}_1' \subset \mathfrak{a}_2' \subset \ldots \subset \mathfrak{a}_m' \neq \mathbf{A}\) ，则 \(m = n\) ，且 \(\mathfrak{a}_k = \mathfrak{a}_k'\) 对 \(1 \leqslant k \leqslant n\) 成立（«诸 \(\mathfrak{a}_k\) 的唯一性»).
